\documentclass[letterpaper]{article} % DO NOT CHANGE THIS
\usepackage[submission]{aaai2027}     % DO NOT CHANGE THIS
% Fonts (newtxtext for serif, helvet for sans, courier for mono) are
% loaded AUTOMATICALLY by aaai2027.sty. Do NOT add \usepackage{times},
% \usepackage{helvet}, or \usepackage{courier} -- they are forbidden in
% the 2027 kit.
\usepackage[hyphens]{url}             % DO NOT CHANGE THIS
\usepackage{graphicx}                 % DO NOT CHANGE THIS
\urlstyle{rm}                         % DO NOT CHANGE THIS
\def\UrlFont{\rm}                     % DO NOT CHANGE THIS
\usepackage{natbib}                   % DO NOT CHANGE THIS AND DO NOT ADD ANY OPTIONS TO IT
\usepackage{caption}                  % DO NOT CHANGE THIS AND DO NOT ADD ANY OPTIONS TO IT
\frenchspacing                        % DO NOT CHANGE THIS
\pdfinfo{
/TemplateVersion (2027.1)
}
% IMPORTANT BEFORE SUBMISSION (anonymous double-blind):
%   AAAI-27 requires the final PDF's metadata be cleared with a
%   metadata-cleaning tool (e.g. exiftool) so authoring info doesn't
%   leak via PDF metadata fields. Do this AFTER the final compile,
%   right before uploading. Example:
%     exiftool -all:all="" -overwrite_original main.pdf
\usepackage{amsmath}
\usepackage{amsthm}
\usepackage{amssymb}
\usepackage{booktabs}
\usepackage{xcolor}                   % required by ReproducibilityChecklist.tex

% --- Revision helper: visible placeholder for content only the authors can
% --- supply. STRIP BEFORE SUBMISSION: grep for \todofill; none may remain.
\newcommand{\todofill}[1]{\textcolor{red}{\textbf{[TO FILL: #1]}}}

% Numbered sections + subsections (theorems counted by section)
\setcounter{secnumdepth}{2}
% Theorem environments
\newtheorem{theorem}{Theorem}[section]
\newtheorem{proposition}[theorem]{Proposition}
\newtheorem{lemma}[theorem]{Lemma}
\newtheorem{observation}[theorem]{Observation}
\newtheorem{corollary}[theorem]{Corollary}
\theoremstyle{definition}
\newtheorem{definition}[theorem]{Definition}
\theoremstyle{remark}
\newtheorem{remark}[theorem]{Remark}
\title{How Does Function Splitting Improve Belief Propagation?}
\author{
    Anonymous Submission
}
\affiliations{
}
\begin{document}
\maketitle
% ========== BEGIN sec0_abstract.tex ==========
\begin{abstract}
Min-sum is the version of belief propagation that is used for solving constraint optimization problems (COPs). Belief propagation is a message passing algorithm in which nodes of the constraint graph (a.k.a.\ factor-graph) take an active role in the algorithm, i.e., they perform calculations and exchange messages with their neighbors. Thus, it is a leading candidate for solving distributed COPs (DCOPs). Unfortunately, the standard vanilla version of Min-sum fails to converge and produces low quality solutions. To mitigate this, previous studies proposed the use of message damping, which encourages the convergence of the algorithm and improves its solution quality. However, damped Min-sum (DMS) requires thousands of iterations in order to converge, and therefore it cannot be used when rapid decision making is required.
Recently, empirical evidence indicated that by splitting the functions of the factor graph on which DMS performs, the algorithm converges immensely faster (almost instantly) to high quality solutions. However, while this success was empirically validated, a theoretical understanding of this phenomenon has not yet been established.

We present theoretical and empirical results that shed light on the success of function-node splitting in improving DMS. We prove that the small cycles generated by such splitting create feedback loops that rapidly separate the solution to which the algorithm converges from alternative candidates, and we establish explicit bounds within which the variable assignments cannot change, i.e., within which the algorithm converges --- including under bounded, dynamically changing influence from the rest of the factor graph. This understanding allows us to propose new heuristics that further improve the algorithm on structured benchmarks, where they significantly outperform a state-of-the-art learned belief-propagation solver.
\end{abstract}
% AAAI-27 \links block: must sit BETWEEN abstract and main body.
% TODO(authors): replace placeholder with the actual anonymised repo URL
% (code + Lean proofs) before submission, and make sure the repo is
% populated. AAAI warns: do not de-anonymise yourself via these links.
\begin{links}
    \link{Code and Lean proofs}{https://anonymized.example/PLACEHOLDER}
\end{links}

% ========== END sec0_abstract.tex ==========

% ========== BEGIN sec1_introduction.tex ==========
\section{Introduction}\label{sec:intro}
Constraint optimization is a general model for centralized and distributed problem solving, which has a wide range of applications in artificial intelligence and multiagent systems \cite{YedidsionZF18,BoothB19,GolestanianBTB23}.
Among the applications that can be represented as constraint optimization problems (COPs) are meeting scheduling, task allocation, bus driver scheduling, disaster response, and operating room assignment \cite{csplib,NelkeOZ20,LavieCLZ23,GaonSZ23}.
Two main approaches are commonly used to solve COPs. The first is search, in which a candidate solution is maintained and updated while traversing the solution space to find one with higher quality than those found previously \cite{ModiSTY05,Zivan08,YeohFK10,LeeMS15}.
	The second is dynamic programming inference \cite{PetcuF05,RogersFSJ11,CohenGZ20} in which information is aggregated so that it enables the selection of high-quality solutions. These inference algorithms are rooted in techniques such as belief propagation and bucket elimination \cite{Pearl88,Dechter99}, which are used to solve optimization problems represented by graphical models.

	Min-sum (also called Max-sum when applied to maximization problems) is an incomplete inference algorithm that has drawn considerable attention in recent years \cite{FarinelliRPJ08,YuCHD17,CerquidesEPR18,DengA20}. It is designed as a message passing algorithm in which nodes in a factor graph exchange messages with neighboring nodes. Thus, it was found to be useful for solving multi-agent applications such as sensor systems \cite{TeacyFGPRJ08,StrandersFRJ09,YedidsionZF18}, task allocation for rescue teams in disaster areas \cite{RamchurnFMJ10}, and IoT applications \cite{RustPR16}.  Min-sum is actually a version of the well known belief propagation algorithm \cite{Pearl88,YanoverMW06}, adjusted to solve COPs and distributed COPs (DCOPs) \cite{BritoM10,GhoshKV15,KhanTYJ18,CohenLZ23}.

	Belief propagation in general (and Min-sum specifically) is known to converge to the optimal solution when solving problems whose constraint graph is acyclic. For problems that include cycles, there is no such guarantee \cite{YanoverMW06,FarinelliRPJ08}.  Furthermore, in graphs with multiple cycles, duplicated information is propagated, leading to inaccurate and inconsistent inference \cite{Pearl88}.  Unfortunately, the underlying problem representation graphs of many realistic applications indeed include multiple cycles (e.g., \citet{ModiSTY05,GershmanMZ09}). Damping is a method that improves Min-sum's performance on problems with multiple cycles~\cite{CohenGZ20,ZaedLZ25}. However, while it triggers convergence of the algorithm to high quality solutions, thousands of iterations are required to reach these, and thus, damped Min-sum (DMS) is not a good choice when very fast decision making is required.
In \citet{CohenGZ20}, the splitting of nodes in the factor-graph upon which the algorithm performs, which represent functions (i.e., function-nodes), was found to trigger rapid convergence of DMS. The DMS algorithm was performed on a new generated factor-graph in which each of the function-nodes in the original factor-graph was represented by two function-nodes. The costs for combination of assignments in the original function-nodes were divided between the two function-nodes in the new generated factor-graph. This caused rapid convergence to high quality solutions. However, although the success of this version of the algorithm (DMS-SCFG) was demonstrated empirically, a formal theoretic explanation of this phenomenon was not provided.

In this study, we investigate the success of function-node splitting in triggering convergence of Min-sum in general and of DMS in particular. We identify a phenomenon that, to the best of our knowledge, has not been reported before: the small cycles generated by the splitting trigger a feedback loop that rapidly separates the solution to which the algorithm converges from alternative candidates. We establish explicit bounds within which the assignments of the variables cannot change, i.e., within which the algorithm converges, and we prove that this convergence survives bounded, dynamically changing influence from the rest of the factor graph.

We further reveal why the use of splitting without damping commonly results in oscillations between solutions of low quality. We demonstrate that Min-sum, when performed on a split constraint factor graph with no damping, tends to converge to two separate solutions between which the nodes of the graph alternate, unfortunately in an uncoordinated manner. This generates mixed assignments of low quality, between which the algorithm oscillates. This explains the importance of combining splitting with damping: damping enforces a strong relation between the calculations performed by nodes in consecutive iterations and thus prevents this phenomenon.

Finally, we use the theoretical knowledge we acquired to design novel heuristics that improve the performance of the algorithm on most benchmarks. The first lets standard Min-sum with splitting reach the stage where it oscillates between two solutions, and then uses standard search techniques to select the best solution among the two. The second runs DMS for a predetermined number of iterations and only then splits the factor graph.

% ========== END sec1_introduction.tex ==========

% ========== BEGIN sec2_related_work.tex ==========
\section{Related Work}
\label{Sec:related}
As mentioned above, \citet{CohenGZ20} demonstrated that function-node splitting triggers rapid convergence of DMS. When applied to standard MS without damping, in most cases it alternates between two solutions of low quality. The DMS-SCFG version converges quickly to a solution of much higher quality. This unexplained success motivated this study.

A general formulation of splitting in the context of belief-propagation was developed by \citet{ruozzi2013}, which showed convergence to local-minima, and some conditions for it to be a global minimum. This was further developed \cite{rebeschini2017} in the specific domain of consensus problems (i.e., as many agents as possible should agree on the same value), in which under several conditions, min-sum using splitting converges to the global optimum in sub-linear time. While formally impressive, these works do not enlighten our understanding of why splitting improves the chances of finding a good outcome in general domains (compared with regular min-sum). Our work, in contrast, wishes to study and understand the \emph{dynamics} of the min-sum mechanism, and how splitting changes these dynamics, ending up with better outcomes.

\citet{DengKL022} identified that the DMS version that uses splitting employs a constant damping parameter for the entire graph. Thus, they proposed an attentive version of the algorithm, in which a deep neural network selects independent damping parameters for every edge of the factor-graph, together with an additional parameter that reduces the weight of incoming messages. With the addition of a specific type of constant splitting, they dramatically improved the quality of solutions produced by DMS-SCFG on benchmarks with random constraints.
This version is not compatible with distributed execution, since the attentive network is trained and applied centrally, with access to the messages of the entire factor graph.% AUTHORS: please verify this characterization of DABP's centralization before submission.
Moreover, as we demonstrate, it loses its advantage on benchmarks with structured constraints. Nevertheless, its success on benchmarks with random constraint costs is striking, and in future work we intend to investigate the patterns of parameter selection that lead to this success and to provide explanations for them.
Most importantly, we demonstrate in our empirical study that the performance of Attentive DMS strongly depends on its use of splitting: without it, the algorithm performs similarly to standard DMS. This is another indication of the importance of investigating the effect of function-node splitting on the performance of DMS, i.e., it further motivates our work.

Our analysis of the two-solution oscillation (Section~\ref{sec:tworoutes}) connects Min-sum on SCFGs to a classical line of results on parallel dynamics with symmetric interactions: synchronous iteration of threshold networks with symmetric weights reaches only fixed points and two-cycles \citep{GolesOlivos80}, as does synchronous weighted-plurality opinion dynamics over finite opinion sets \citep{PoljakSura83}; a period-two result for simultaneous best-response dynamics was recently established for two-player random potential games \citep{AshkenaziGolan25}. These results concern special interaction forms (thresholds, equality indicators, two players); the reduction of committed split Min-sum to synchronous local selection, and the resulting two-solution characterization for arbitrary shared pairwise cost terms, are, to the best of our knowledge, new.

% ========== END sec2_related_work.tex ==========

% ========== BEGIN sec3_background.tex ==========
\section{Background}
\label{Sec:background}
In this section, we present the preliminaries of distributed constraint optimization problems (DCOPs) and the version of belief propagation used to solve them -- the min-sum algorithm and its version that includes damping.
\subsection{Distributed Constraint Optimization Problems (DCOPs)}
\label{Sec:DCOP_defs}
\sloppy
A DCOP\footnote{For lack of space and to avoid redundancy, we do not provide a separate description of a centralized COP. A COP can be considered a DCOP with a single agent.} is a tuple $\langle{\mathcal A},{\mathcal X},{\mathcal D},{\mathcal C}\rangle$, where ${\mathcal A}$ is a finite set of agents $\{ A_1,A_2,\ldots,A_n \}$; ${\mathcal X}$ is a finite set of variables $\{ X_1,X_2,\ldots,X_m \}$, where each variable is held by a single agent (an agent may hold more than one variable); ${\mathcal D}$ is a set of domains $\{ D_1, D_2,\ldots,D_m \}$, where each domain $D_i$ contains the finite set of values that can be assigned to variable $X_i$ and we denote an assignment of value $d \in D_i$ to $X_i$ by an ordered pair $\langle X_i, d \rangle$; and ${\mathcal C}$ is a set of constraints (relations), where each constraint $C_j \in {\mathcal C}$ defines a non-negative {\em cost} for every possible value combination of a set of variables and is of the form $C_j : D_{j_1} \times D_{j_2} \times \ldots \times D_{j_k} \rightarrow \mathbb{R}^+ \cup \{0\}$. A {\em binary constraint} refers to exactly two variables and is of the form $C_{ij} : D_i \times D_j \rightarrow \mathbb{R}^+ \cup \{0\}$.
In our discussion of Min-sum we often refer to $C_{ij}$ as the {\em cost table} and to $C_{ij}[w,r]$ as the entry corresponding to $w \in D_i$ and $r \in D_j$.

A {\em binary DCOP} is a DCOP in which all constraints are binary. Agents are {\em neighbors} if they are involved in the same constraint. A {\em partial assignment} (PA) is a set of value assignments to variables, in which each variable appears at most once. A constraint $C_j \in {\mathcal C}$ of the form $C_j : D_{j_1} \times D_{j_2} \times \ldots \times D_{j_k} \rightarrow \mathbb{R}^+ \cup \{0\}$ is {\em applicable} to $PA$ if each of the variables $X_{j_1} , \ldots , X_{j_k}$ is included in PA. The {\em cost of a partial assignment} $PA$ is the sum of all applicable constraints to $PA$ over the value assignments in $PA$. A {\em complete assignment} (i.e., {\em solution}) is a partial assignment that includes all variables ($\mathcal X$). An {\em optimal solution} is a complete assignment with minimal cost.
For simplicity, we assume that each agent holds exactly one variable (i.e., $n=m$) and focus on binary DCOPs. These assumptions are common in DCOP literature (e.g.,~\cite{Petcu05b,YeohFK10}).
Min-sum's complexity is exponential in the arity of constraints, and therefore its applicability to problems that include constraints with high arity is limited~\cite{FarinelliRPJ08,YedidsionZF18}. However, this limitation is not relaxed by the use of damping, therefore, it is out of the scope of this study. We note that in contrast to Attentive-DMS, the versions of Min-sum we propose in this study do apply to factor-graphs with non binary function-nodes.
\subsection{Min-sum Belief Propagation}
\label{secmaxsum}
Min-sum operates on a \emph{factor-graph}, a bipartite graph in which the nodes represent variables and functions (constraints)~\cite{KschischangFL01}.\footnote{We preserve the terminology of~\cite{FarinelliRPJ08} and call constraint representing nodes in the factor-graph ``function-nodes".}
Each variable-node, representing a variable of the original DCOP, is connected to all function-nodes representing constraints that it is involved in. In other words, a function-node is connected to all variable-nodes involved in the constraint it represents.
Variable-nodes and function-nodes take an active role in Min-sum, i.e., they can send and receive messages, and can perform computation. When used to solve a DCOP, each node's role is performed by an autonomous agent.
In Min-sum, a message sent to -- or from -- variable-node $X$ (for simplicity, we use the same notation for a variable and the variable-node representing it) is a vector of size $|D_X|$ including a cost (or the belief of a cost) for each value in $D_X$.
In the first iteration, all nodes assume that all messages they previously received (in iteration $0$) include vectors of zeros. Each iteration is \emph{two-phase}: first all variable-nodes compute and send their messages, then all function-nodes compute and send theirs using the variable messages of the same iteration (this is the schedule implemented and analyzed throughout the paper). A message $Q_{X \rightarrow F}^i$ that variable-node $X$ sends to function-node $F$ in iteration $i$ is formalized as follows:
$Q_{X \rightarrow F}^i =\underset{F'\in F_X \setminus \{ F\} }{ \sum} R_{F' \rightarrow X}^{i-1} - \alpha$,
where $F_X$ is the set of function-node neighbors of variable-node $X$ and $R_{F' \rightarrow X}^{i-1}$ is the message sent to variable-node $X$ by function-node $F'$ in iteration $i-1$.
$\alpha$ is a constant that is reduced from all costs included in the message in order to prevent the costs from growing arbitrarily large.
A message $R_{F \rightarrow X}^{i}$ sent from a function-node $F$ to a variable-node $X$ in iteration $i$,
includes for each value $x \in D_X$:
$R_{F \rightarrow X}^{i} = \underset{{PA_{-X}}} \min\, cost (\langle X , x \rangle, PA_{-X})$,
where $PA_{-X}$ is a possible combination of value assignments to variables involved in $F$ not including $X$. $cost(\langle X , x \rangle, PA_{-X})$ represents the cost of a partial assignment $a = \{ \langle X, x \rangle, PA_{-X} \}$, which is:
$f(a) + \underset{X' \in X_F \setminus \{X\},\langle X',x' \rangle \in a}{\sum} (Q_{X' \rightarrow F}^{i})[x']$,
where $f(a)$ is the original cost in the constraint represented by $F$ for the partial assignment $a$, $X_F$ is the set of variable-node neighbors of $F$, and $(Q_{X' \rightarrow F}^{i})[x']$ is the cost in the message sent from variable-node $X'$ in the first phase of the same iteration, for the value $x'$ that is assigned to $X'$ in $a$.
$X$ selects its value assignment $\hat{x} \in D_X$ following iteration $k$ as follows:
$\hat{x} = \underset{x\in D_X}{\mathrm{argmin}} \sum_{F\in F_X} (R_{F \rightarrow X}^k)[x]$.
\subsection{Damped Min-sum (DMS)}
In order to add damping to Min-sum, a parameter $\lambda \in [0,1)$ is used. Before sending a message in iteration $k$, a node performs calculations as in standard Min-sum. We use $\widehat{m^k_{i \rightarrow j}}$ to denote the result of the calculation made in standard Min-sum for the content of a message intended to be sent from node $i$ to node $j$ in iteration $k$ and $m^{k-1}_{i \rightarrow j}$ to denote the message sent by $i$ to $j$ at iteration $k-1$. The message sent by $i$ to $j$ at iteration $k$ is calculated as follows:
$m^{k}_{i \rightarrow j} = \lambda m^{k-1}_{i \rightarrow j} + (1 - \lambda) \widehat{m^k_{i \rightarrow j}}$.
Thus, $\lambda$ expresses the weight given to previous calculations performed with respect to the most recent calculation performed. When $\lambda = 0$, the resulting algorithm is standard Min-sum.
\subsection{Backtrack Cost Trees}
\label{sec:bct}
	For analyzing the behavior of Min-sum, the use of a {\em backtrack cost tree} (BCT) was proposed by \citet{ZivanLG20}. It allows tracing for each belief the entries in the cost tables held by function-nodes that were used to compose it. In other words, the components of the assignment's cost.
A BCT is defined for a belief that appears in a message (either from some variable $X_{i}$ to a function-node $F_{ij}$, or from some function-node $F_{ij}$ to a variable $X_{i}$). The belief is on the cost of assigning some value $x\in D_{i}$ to variable $X_{i}$. Without loss of generality, we will elaborate on messages from variables to function-nodes.
The belief is a sum of various components from which the BCT is composed. The root is the cost of a decision to assign some value to a variable at a time $t$ and the directed edges from its children in the tree include the beliefs that were summed to generate that belief. These edges lead to nodes representing the neighbors from which the parent node received messages in time $t-1$. Each of those nodes is connected to the nodes from which it received messages at time $t-2$, with the edges containing the beliefs that were passed to it. The tree leaves represent the nodes at time $0$.

% ========== END sec3_background.tex ==========

% ========== BEGIN sec4_effect_of_splitting.tex ==========
% ==================================================================
\section{Effect of Splitting}\label{sec:setting}
% ==================================================================
In this section, we present our key insight --- an analysis of the effect of splitting on the behavior of the Min-sum algorithm and the reason it triggers convergence.

\paragraph{Scope of the analysis.}
% NEW in revision: explicit scope statement (reviewer-proofing).
The formal results in this section concern undamped Min-sum on a single split binary function-node with two-value domains: Sections~\ref{sec:regimes}--\ref{sec:convergence} treat the isolated split cycle, and Section~\ref{sec:dynamic} models the influence of the remainder of the factor graph as a bounded, dynamically changing unary constraint. Section~\ref{sec:extensions} takes first steps toward damped updates and asymmetric splits. The role of damping in full graphs is analyzed in Sections~\ref{Sec:MS-split}--\ref{sec:whydamping}, and the instrumentation experiment of Section~\ref{sec:mechanism} measures how often the assumptions made here hold on benchmark instances.
A note on indexing: within this section, message superscripts stamp the successive messages traveling \emph{around the split cycle} (so that consecutive messages of the same type are two stamps apart, e.g., $\Delta Q^{t-2}$ determines $\Delta R^{t}$); this differs from the two-phase iteration count of Section~\ref{secmaxsum} only by a relabeling of the time axis, and none of the results depends on the absolute stamps --- only on the traversal order of the cycle.

\begin{figure}
\centering
\includegraphics[width=8cm]{examp.pdf}
\vspace{-10pt}
\caption{A factor graph with a single binary function-node converted into a split constraint factor graph with a single cycle.}
\vspace{-5pt}
\label{Fig:examp}
\end{figure}
We start by discussing the effect of splitting on the most simple case of a single binary function-node $F_{12}$ with two variable-nodes $X_1, X_2$ with two values in each of their domains  $D_i = \{a, b\}$ as depicted in Figure~\ref{Fig:examp}. Before starting the algorithm iterations, we split $F_{12}$ to $F'_{12}$ and $F^{''}_{12}$ evenly, i.e., for every entry $e$ in the cost table of $F_{12}$, $e' = e^{''} = \frac{1}{2}e$ (where $e'$ and $e^{''}$ are the entries that correspond to entry $e$ in $F'_{12}$ and $F^{''}_{12}$, respectively).
To make the presentation more simple, we gave the following notations to the entries: $M_a$ to $F_{12}[a,a]$,  $B_a$ to $F_{12}[b,a]$, $B_b$ to $F_{12}[a,b]$, and $M_b$ to $F_{12}[b,b]$ (see Figure~\ref{Fig:examp}). We will also assume w.l.o.g.\ that $M_a < M_b < B_b < B_a$; intuitively, one may think of both $B_a$ and $B_b$ as much larger than $M_a$ and $M_b$, but only the strict order is used in the formal results.
We further denote by $Q^t_{i} = \langle Q^t_{ia},Q^t_{ib} \rangle$ the message sent from $X_i$ to $F_{12}$ at iteration $t$ and by $R^t_{i} = \langle R^t_{ia},R^t_{ib} \rangle$ the message sent from $F'_{12}$ to $X_i$ at iteration $t$.
We denote by $\Delta_Q$ the difference between the entries, i.e., $\Delta_Q = Q_b - Q_a$, the margin by which the value with the smaller entry is preferred. Moreover, the minimal entry in a message is the smallest among $Q_a$ and $Q_b$.
Similar notation will be used for messages $R$.
    For a message $R = \langle R_a, R_b \rangle$, $\Delta_R = R_b - R_a$ is the margin by which the value with the smaller entry is preferred. A variable-node $X$ changes the value implied by $R$ only if the differences it receives from its other neighbors sum, in the opposite direction, to more than $\Delta_R$, i.e., when $\underset{R' \neq R}{\sum}\Delta_{R'} < -\Delta_R$.
\begin{observation}\label{obs:unsplit}
In the original (un-split) factor graph, $X_2$ has a single neighbor, so $R^t_{2} = \langle M_a, M_b \rangle$ is static: $\Delta_{R^t_2} = \Delta_{R^{t'}_2} = M_b - M_a$ for all iterations $t, t'$.
\end{observation}
Recall that to simplify the discussion, we assumed that in our example $M_a < M_b$ and therefore in the original factor graph, $\Delta_R = M_b - M_a$ at all iterations.
However, after the split, a feedback loop is generated that updates the content of the messages and specifically $\Delta_R$.
In more detail, this is the sequence of messages that are passed in the first few iterations of the algorithm when solving the single cycle factor graph on the right hand side of Figure~\ref{Fig:examp}:\footnote{Since following the first iteration all messages $Q$ are equal to the messages $R$ received by the variable-node sending them in the previous iteration, we only specify the content of the $R$ messages.} $R^1_{2'} = \langle M'_a,M'_b\rangle$, $R^3_{1''} = \langle 2M'_a,2M'_b\rangle$, $R^5_{2'} = \langle 3M'_a,3M'_b\rangle$,  $R^7_{1''} = \langle 4M'_a,4M'_b\rangle$...
Thus, $\Delta_Q$ and $\Delta_R$ grow by $M'_b - M'_a$ every two iterations. This process will continue until $\Delta_R = B'_b - M'_a$. Then, the algorithm will converge and the same messages will be sent by the nodes in the graph in all following iterations. We formally prove this in the following Lemma:
\begin{lemma}\label{lem:singlecycle}
    In an SCFG with a single cycle in which the entries are denoted as depicted in Figure~\ref{Fig:examp} and w.l.o.g. $M_a < M_b < B_b < B_a$, $\Delta_{R_1}$ will converge to $B'_b - M'_a$.
\end{lemma}
\begin{proof}
    A variable-node in the single cycle forwards the message it received from its other neighbor, hence $\Delta^{t}_Q = \Delta^{t-1}_R$, starting from $\Delta^0_Q = 0$. As long as $M'_b + \Delta^t_Q < B'_b$, both entries of $R$ take the ``$M$'' branch, so $R_a = M'_a + Q_a$ and $R_b = M'_b + Q_b$, giving $\Delta^{t+1}_R = (M'_b - M'_a) + \Delta^t_Q$; thus $\Delta_R$ grows by $M'_b - M'_a$ every two iterations. This continues until $\Delta^t_Q \geq B'_b - M'_b$, where the $b$-entry switches to the ``$B'_b$'' branch, $R_b = B'_b + Q_a$, giving $\Delta_R = (B'_b + Q_a) - (M'_a + Q_a) = B'_b - M'_a$, independent of $Q$. Substituting back, $\Delta_Q = B'_b - M'_a$ keeps the $B'_b$ branch active (since $M'_b \geq M'_a$) and keeps $R_a$ on the $M'_a$ branch (since $2M'_a < B'_a + B'_b$), so $\Delta_R$ remains $B'_b - M'_a$: the process has converged.
\end{proof}
It is important to note that if $B_a < B_b$, then the algorithm would converge when $\Delta_R = B'_a - M'_a$. Thus, we will call $B_a$ and $B_b$ the bounders.
This is the first indication of the implications of the split of a function-node. The increase of $\Delta_R$ results in a much more stable state, since the $\Delta$ of the messages that the variable-nodes $X_1$ and $X_2$ must receive to change the assignment selection implied by the $R$ messages received from $F'_{12}$ and $F^{''}_{12}$ is much larger. In other words, the algorithm generates larger differences between its minimal selection of assignments and the cost of the alternatives, triggering convergence.
Next, we will analyze an extended scenario in which the variable-nodes in the cycle receive messages from function-nodes that are not part of the cycle.
For each value $v$ of $X_2$, the forward update compares the two assignments of $X_1$: it adds $Q_a$ to the entry with $X_1=a$ and $Q_b$ to the entry with $X_1=b$, and keeps the smaller.  The incoming message thus shifts the cost of each assignment of $X_1$, and the minimum selects the cheaper one for every value of $X_2$.  This is what the forward min-sum update of Section~\ref{secmaxsum}, specialised to $C'$, says formally:
\begin{equation}\label{eq:minsumR}
    \bigl(R^t_{F_{12} \to X_2}\bigr)_v = \min_{u \in D}\bigl[C'(u, v) + (Q^{t-2}_{X_1 \to F_{12}})_u\bigr], \;\; v \in D.
\end{equation}
Throughout this section $C'$ denotes the split copy of the function-node, and we write $M'_a := C'(a,a)$, $B'_a := C'(b,a)$, $B'_b := C'(a,b)$, $M'_b := C'(b,b)$ for its four entries (Figure~\ref{Fig:examp}); the results below hold for arbitrary such entries. Which assignment of $X_1$ is selected for each value of $X_2$ depends only on the \emph{gap} between the two shifted entries; the absolute level of $Q$ cancels, so the analysis reduces to the message difference.
\begin{definition}[Message differences]\label{def:deltas}
Let
\[
    \begin{aligned}
        \Delta Q^t &:= \bigl(Q^t_{X_1 \to F_{12}}\bigr)_b - \bigl(Q^t_{X_1 \to F_{12}}\bigr)_a, \\
        \Delta R^t &:= \bigl(R^t_{F_{12} \to X_2}\bigr)_b - \bigl(R^t_{F_{12} \to X_2}\bigr)_a.
    \end{aligned}
\]
\end{definition}
\noindent The recurrence is initialised by an arbitrary real-valued message vector; equivalently, we treat the initial difference $\Delta Q^0 \in \mathbb{R}$ as a free parameter.
As $\Delta Q$ varies, the forward step's outcome goes through three qualitatively different cases: for $\Delta Q$ very positive, $X_1=a$ is selected for both values of $X_2$; in a middle range the two values of $X_2$ select different assignments; for $\Delta Q$ very negative, $X_1=b$ is selected for both.  The switch happens separately for each value of $X_2$: each $v$ has a critical $\Delta Q$ at which its two shifted entries tie and the selected assignment of $X_1$ flips.
\begin{definition}[Assignment margins and stability thresholds]\label{def:thresholds}
For each value $v \in D$ of $X_2$ define the margin of assignment $X_1=a$ over $X_1=b$,
\[
    \Delta_v \;:=\; C'(X_1{=}a,\, X_2{=}v) - C'(X_1{=}b,\, X_2{=}v),
\]
so that $X_1=a$ is selected for $X_2=v$ iff $\Delta Q \geq \Delta_v$. Without loss of generality (by permuting the labels $\{a,b\}$ of $X_2$ if necessary) assume $\Delta_a \leq \Delta_b$, and set the \emph{stability thresholds}
\[
    \tau_U \;:=\; \Delta_b, \qquad \tau_L \;:=\; \Delta_a,
\]
so that $\tau_L \leq \tau_U$.
\end{definition}
\noindent Concretely $\Delta_a = M'_a - B'_a$ and $\Delta_b = B'_b - M'_b$; for the agreement structure $M'_a < M'_b < B'_b < B'_a$ of Section~\ref{sec:setting} this gives $\tau_U = B'_b - M'_b > 0$ and $\tau_L = M'_a - B'_a < 0$: for $\Delta Q > \tau_U$ the message selects $X_1=a$ for both values of $X_2$, for $\Delta Q < \tau_L$ it selects $X_1=b$, and in between the two values of $X_2$ select different assignments.  Outside the interval $[\tau_L, \tau_U]$ the forward step is constant in $\Delta Q$; inside, the selection \emph{splits}.  The remainder of this section formalizes what these regimes mean for the message dynamics.
% ------------------------------------------------------------------
\subsection{Regime Stability}\label{sec:regimes}
% ------------------------------------------------------------------
When $\Delta Q$ falls outside the interval $(\tau_L, \tau_U)$, $F_{12}$ selects the same assignment of $X_1$ for every value of $X_2$, and the outgoing difference $\Delta R$ is therefore determined by $C'$ alone, regardless of the incoming message.
\begin{theorem}[Stability of the outgoing message]\label{thm:stability}
Let $\tau_L, \tau_U$ be as in Definition~\ref{def:thresholds}.
\begin{enumerate}
    \item \textbf{(Upper regime.)} If $\Delta Q^{t-2} \geq \tau_U$, then
    \[
        \Delta R^t \;=\; C'(a,b) - C'(a,a) \;=\; B'_b - M'_a,
    \]
    independently of $Q^{t-2}_{X_1 \to F_{12}}$.
    \item \textbf{(Lower regime.)} If $\Delta Q^{t-2} \leq \tau_L$, then
    \[
        \Delta R^t \;=\; C'(b,b) - C'(b,a) \;=\; M'_b - B'_a,
    \]
    independently of $Q^{t-2}_{X_1 \to F_{12}}$.
\end{enumerate}
In both cases $\Delta R^t$ depends only on $C'$.
\end{theorem}
\begin{proof}
By~\eqref{eq:minsumR}, for each value $v$ of $X_2$ the minimum is attained at $X_1=a$ iff $\Delta Q^{t-2} \geq \Delta_v$, and at $X_1=b$ iff $\Delta Q^{t-2} \leq \Delta_v$ (ties give the same value regardless).
\smallskip\noindent\textbf{Case 1} ($\Delta Q^{t-2} \geq \tau_U \geq \tau_L$):  $X_1=a$ is selected for both values of $X_2$, so $(R^t)_v = C'(a,v) + (Q^{t-2})_a$.  The $(Q^{t-2})_a$ terms cancel in the difference: $\Delta R^t = C'(a,b)-C'(a,a)$.
\smallskip\noindent\textbf{Case 2} ($\Delta Q^{t-2} \leq \tau_L \leq \tau_U$):  $X_1=b$ is selected for both values of $X_2$, so $(R^t)_v = C'(b,v) + (Q^{t-2})_b$.  The $(Q^{t-2})_b$ terms cancel: $\Delta R^t = C'(b,b)-C'(b,a)$. \qedhere
\end{proof}
Once $\Delta Q$ is in either regime, the outgoing message is therefore \emph{locked}: $\Delta R$ equals a difference of two entries of $C'$ and is severed from the incoming $Q$ entirely.  We call $(\tau_L, \tau_U)$ the \emph{transition interval}: when $\Delta Q$ lies in this open interval, neither case applies and $\Delta R$ depends on the incoming message.
% ------------------------------------------------------------------
\subsection{Absorbing Regimes}\label{sec:absorbing}
% ------------------------------------------------------------------
Having identified the regimes, we ask: once the process enters a regime, does it stay? A single step into a regime produces a constant $\Delta R$, which feeds back to yield a specific next value of $\Delta Q$—the \emph{round-trip constant}. The regime is absorbing precisely when this constant remains inside the regime.
The backward pass is the mirror image of~\eqref{eq:minsumR}: instead of fixing a value of $X_2$ and minimising over the assignments of $X_1$, we fix a value of $X_1$ and minimise over the assignments of $X_2$.  Using the same split cost $C'$, the updated message from $X_1$ to $F_{12}$ is
\begin{equation}\label{eq:minsumQ}
    \bigl(Q^t_{X_1 \to F_{12}}\bigr)_{x_1} = \min_{x_2 \in D}\bigl[C'(x_1, x_2) + (R^{t-2}_{F_{12} \to X_2})_{x_2}\bigr].
\end{equation}
Suppose the trajectory has just entered a regime: $\Delta Q^{t-2}$ is at or beyond a threshold.  The forward step then produces the regime's constant $\Delta R^t$ (Theorem~\ref{thm:stability}), and feeding this constant into the backward step yields a specific value of $\Delta Q^{t+2}$, also independent of the message level.  We call these the lower and upper \emph{round-trip constants}; if a constant lies inside its own regime, the trajectory is trapped there for good.  Both observations are packaged into the following theorem.
\begin{theorem}[Round-trip constants and absorption]\label{thm:absorbing}
\begin{enumerate}
\item \textbf{(Lower.)}  If $\Delta Q^{t-2} \leq \tau_L$ for some $t \geq 2$, then $\Delta R^t = C'(b,b) - C'(b,a)$ and
\begin{multline}\label{eq:cL}
    \Delta Q^{t+2}
    = \min\!\bigl(2C'(b,a),\;2C'(b,b)\bigr) \\
    - \min\!\bigl(C'(a,a)+C'(b,a),\;C'(a,b)+C'(b,b)\bigr),
\end{multline}
independently of $Q^{t-2}$ and of $t$; denote this value by $c_L(C')$.  If additionally $c_L \leq \tau_L$, then $\Delta Q^j \leq \tau_L$ and $\Delta R^j = C'(b,b) - C'(b,a)$ for all $j \geq t$ --- the regime is \emph{absorbing}.
\item \textbf{(Upper.)}  If $\Delta Q^{t-2} \geq \tau_U$ for some $t \geq 2$, then $\Delta R^t = C'(a,b) - C'(a,a)$ and
\begin{multline}\label{eq:cU}
    \Delta Q^{t+2}
    = \min\!\bigl(C'(a,a)+C'(b,a), \\
    C'(a,b)+C'(b,b)\bigr)
    - \min\!\bigl(2C'(a,a),\,2C'(a,b)\bigr),
\end{multline}
independently of $Q^{t-2}$ and of $t$; denote this value by $c_U(C')$.  If additionally $c_U \geq \tau_U$, then $\Delta Q^j \geq \tau_U$ and $\Delta R^j = C'(a,b) - C'(a,a)$ for all $j \geq t$ --- the regime is \emph{absorbing}.
\end{enumerate}
\end{theorem}
\begin{proof}
\textbf{Upper.}  By Theorem~\ref{thm:stability} (Case 1), $(R^t)_v = C'(a,v)+(Q^{t-2})_a$, giving $\Delta R^t = C'(a,b)-C'(a,a)$.  Substituting into~\eqref{eq:minsumQ}, the constant shift $(Q^{t-2})_a$ cancels in the difference, yielding~\eqref{eq:cU}.  If additionally $c_U \geq \tau_U$, then $\Delta Q^{t+2} = c_U \geq \tau_U$, so the same argument applies at $t+4$ and inductively for all $j \geq t$.  \textbf{Lower.}  Symmetric: $X_1=b$ throughout, $(Q^{t-2})_b$ cancels, and the same induction. \qedhere
\end{proof}
\noindent For the agreement structure $M'_a < M'_b < B'_b < B'_a$ of Section~\ref{sec:setting} it is the \emph{upper} regime that is absorbing: $c_U \geq \tau_U$ (established in Theorem~\ref{thm:convergence}), consistent with the single-cycle settling difference $\Delta R = B'_b - M'_a$ of Section~\ref{sec:setting}; the lower regime is not absorbing.
% ------------------------------------------------------------------
\subsection{Convergence to an Absorbing Regime}\label{sec:convergence}
% ------------------------------------------------------------------
Theorem~\ref{thm:absorbing} is conditional: it gives absorption \emph{if} the trajectory enters the regime \emph{and} the round-trip constant lies inside it.  We now show the unconditional version: starting from any $\Delta Q^0$, the trajectory \emph{has} to reach the upper regime in finitely many iterations and stay there.
The driving quantity is the \emph{drift constant}
\[
    d \;:=\; C'(b,b) - C'(a,a) \;=\; M'_b - M'_a.
\]
Inside the transition interval the forward step has $\Delta R = \Delta Q + d$ (for $X_2=a$ the minimum is at $X_1=a$, and for $X_2=b$ at $X_1=b$).  When the backward step is likewise split (for $X_1=a$ the minimum is at $X_2=a$, and for $X_1=b$ at $X_2=b$), a direct computation from~\eqref{eq:minsumQ} gives $\Delta Q^{\text{next}} = \Delta R + d$, contributing another $d$ and yielding a net change of $2d$ per round-trip.  A round-trip is one full loop of the split cycle through both copies $F', F''$ --- four iterations --- so $\Delta Q$ advances by $|d|$ every two, as in Section~\ref{sec:setting}.  Since $d \neq 0$, this monotone drift forces $\Delta Q$ out of the transition interval in finite time.
\textbf{Canonical orientation.}  We assume $d > 0$ and $\tau_U + \tau_L \leq 0$ throughout, which lets us track a single regime (the upper one); the $d < 0$ case follows by symmetry.  Two relabellings get there: swapping $\{a, b\}$ in \emph{both} $X_1$ and $X_2$ sends $d \mapsto -d$ (preserving $\Delta_a \leq \Delta_b$), and swapping $X_1 \leftrightarrow X_2$ replaces $\tau_U + \tau_L$ by a quantity summing with the original to $-2d < 0$, so one of the two values is nonpositive.
Inside the transition interval, each round-trip either exits to the upper regime, or increments $\Delta Q$ by at least $2|d|$.
\begin{lemma}[Drift in the transition interval]\label{lem:drift}
Assume $d > 0$ and $\tau_U + \tau_L \leq 0$.  For any $\Delta Q^t \in (\tau_L, \tau_U)$, after one round-trip either
\[
    \Delta Q^{t+4} \;\geq\; \tau_U \qquad\text{(upper regime entered),}
\]
or $\Delta Q^{t+4} \in (\tau_L, \tau_U)$ with $\Delta Q^{t+4} - \Delta Q^t \geq 2|d|$.
\end{lemma}
\begin{proof}
Normalise $R^{t+2}_a = 0$.  In the transition interval, for $X_2=a$ the minimum is at $X_1=a$ and for $X_2=b$ at $X_1=b$, so $\Delta R^{t+2} = \Delta Q^t + d =: r$.
In the Q-update~\eqref{eq:minsumQ}, assignment $X_1=u$ prefers $X_2 = a$ iff $r \geq C'(u,a)-C'(u,b)$.  Since $\tau_L \leq \tau_U$ implies $C'(a,a)-C'(a,b) \leq C'(b,a)-C'(b,b)$, three cases arise:
\textbf{Both assignments prefer $X_2 = a$} ($r \geq C'(b,a)-C'(b,b)$): $\Delta Q^{t+4} = C'(b,a)-C'(a,a) = -\tau_L \geq \tau_U$.
\textbf{Both prefer $X_2 = b$} ($r \leq C'(a,a)-C'(a,b)$): $\Delta Q^{t+4} = C'(b,b)-C'(a,b) = -\tau_U$.  Using $C'(a,a)-C'(a,b) = -d-\tau_U$, the condition $r \leq -d-\tau_U$ gives $\Delta Q^t = r-d \leq -2d-\tau_U$, so the increase $\Delta Q^{t+4}-\Delta Q^t = -\tau_U-\Delta Q^t \geq 2|d|$.
\textbf{Split} (otherwise): $X_1=a$ prefers $X_2=a$ and $X_1=b$ prefers $X_2=b$, giving $\Delta Q^{t+4} = C'(b,b)+r-C'(a,a) = r+d = \Delta Q^t+2d$, an increase of exactly $2|d|$. \qedhere
\end{proof}
% ------------------------------------------------------------------
\subsubsection{Convergence Theorem}
% ------------------------------------------------------------------
\begin{lemma}[Non-absorption of the lower regime]\label{lem:upper-nonabsorb}
Under the WLOG orientation $d > 0$ and $\tau_U + \tau_L < 0$, we have $c_L > \tau_L$.  Consequently, any trajectory with $\Delta Q^{t-2} \leq \tau_L$ satisfies $\Delta Q^{t+2} = c_L > \tau_L$, and the lower regime is non-absorbing.
\end{lemma}
\begin{proof}
By Theorem~\ref{thm:stability}, the lower regime gives $\Delta R^t = \tau_L + d$, independent of $Q^{t-2}$.  The backward update~\eqref{eq:minsumQ} then computes
\[
    \begin{aligned}
        c_L &= \min\bigl(C'(b,a)+R_a,\; C'(b,b)+R_b\bigr) \\
            &\quad - \min\bigl(C'(a,a)+R_a,\; C'(a,b)+R_b\bigr).
    \end{aligned}
\]
Applying the inequality $\min(A,B) - \min(C,D) \geq \min(A-C,\, B-D)$,
\[
    \begin{aligned}
        c_L &\geq \min\bigl(C'(b,a) - C'(a,a),\; C'(b,b) - C'(a,b)\bigr) \\
            &= \min(-\tau_L,\, -\tau_U) = -\tau_U.
    \end{aligned}
\]
Since $\tau_U + \tau_L < 0$, $-\tau_U > \tau_L$, hence $c_L > \tau_L$. \qedhere
\end{proof}
\begin{theorem}[Convergence to an absorbing regime]\label{thm:convergence}
Assume $d \neq 0$ and let the WLOG relabelling above be in effect.  Starting from any initial $\Delta Q^0 \in \mathbb{R}$, there exists a finite index $N$ such that for all $j \geq N$:
\[
    \Delta Q^j \geq \tau_U,\quad \Delta R^j = C'(a,b) - C'(a,a) \quad (d > 0),
\]
and symmetrically for $d < 0$ with the lower regime.
\end{theorem}
\begin{proof}
We treat $d > 0$.
\textbf{Step 1: reaching the upper regime.}  If $\Delta Q^0 \geq \tau_U$, we are done.  If $\Delta Q^0 \leq \tau_L$, one round-trip gives $\Delta Q^4 = c_L$; by Lemma~\ref{lem:upper-nonabsorb}, $c_L > \tau_L$, so $\Delta Q^4$ either lies in the transition interval or is already in the upper regime.
For iterates in the transition interval, Lemma~\ref{lem:drift} guarantees that each round-trip either enters the upper regime directly, or increases $\Delta Q$ by at least $2|d|$.  Since the transition interval has finite width $\tau_U - \tau_L$, the upper regime is reached in at most $\lceil(\tau_U - \tau_L)/2|d|\rceil$ round-trips.
\textbf{Step 2: the upper regime is absorbing.}  We compute $c_U$ in closed form and verify $c_U \geq \tau_U$.  Write $p = C'(a,a)$, $q = C'(a,b)$, so that $C'(b,a) = p-\tau_L$, $C'(b,b) = q-\tau_U$, and $q-p = d+\tau_U$.  Substituting into~\eqref{eq:cU},
\[
    c_U \;=\; \min(2p-\tau_L,\, 2q-\tau_U) \;-\; \min(2p,\, 2q).
\]
The second minimum splits on the sign of $p-q = -(d+\tau_U)$, equivalently $\tau_U \gtrless -|d|$; the first on the sign of $(2p-\tau_L)-(2q-\tau_U) = -2d - (\tau_U+\tau_L)$, equivalently $\tau_U+\tau_L \gtrless -2|d|$.  A direct case analysis yields
\[
    c_U \;=\;
    \begin{cases}
        -\tau_U, & \tau_U \leq -|d|,\\[2pt]
        \tau_U + 2|d|, & \tau_U > -|d|,\ \tau_U+\tau_L < -2|d|,\\[2pt]
        -\tau_L, & \text{otherwise.}
    \end{cases}
\]
In each case $c_U \geq \tau_U$: the first because $\tau_U \leq -|d| < 0$ gives $-\tau_U \geq |d| \geq \tau_U$; the second because $|d|>0$; the third because $\tau_U+\tau_L \leq 0$ gives $-\tau_L \geq \tau_U$.  Theorem~\ref{thm:absorbing} (Part 2) therefore applies.  The margin satisfies
\[
    c_U - \tau_U \;\geq\; \min\!\bigl(2|d|,\; -(\tau_U+\tau_L)\bigr),
\]
with equality when $\tau_U \geq -|d|$; for $\tau_U < -|d|$ the margin equals $-2\tau_U$, strictly larger.  \qedhere
\end{proof}
This completes the static picture: splitting induces a regime structure on the message difference, with the gating thresholds $\tau_L, \tau_U$ read off $C'$; the trajectory is forced into the upper regime in finite time from any start, where it settles at $\Delta R = B'_b - M'_a$.
% ------------------------------------------------------------------
\subsection{Convergence under a Dynamic Unary Constraint on $X_1$}\label{sec:dynamic}
% ------------------------------------------------------------------
In the static setting of Section~\ref{sec:setting} the unary constraint on $X_1$ was fixed.  We now let it vary.  When $X_1$ has other neighbours in a larger factor graph, the message it sends to $F_{12}$ carries, at each iteration $t$, a unary constraint $\phi^t : D \to \mathbb{R}$ that changes from step to step.
Each iteration folds in a \emph{different} $\phi^t$, giving a different effective cost $C'^{(t)}$ and hence different gating thresholds.  The static analysis still carries over while the drift keeps pushing $\Delta Q$ toward a single regime; near a regime boundary, however, the shifting thresholds can change which assignment of $X_1$ the message selects.  We now characterise when finite-time convergence survives.
\begin{definition}[Effective cost at iteration $t$]\label{def:dynC}
\[
    \begin{aligned}
        C'^{(t)}(x_1, x_2) &:= C'(x_1, x_2) + \phi^t(x_1), \\
        \delta^t_\phi &:= \phi^t(b) - \phi^t(a).
    \end{aligned}
\]
\end{definition}
\begin{definition}[Dynamic thresholds]\label{def:dynthresh}
For each $v \in D$ set $\Delta^t_v := C'^{(t)}(a,v) - C'^{(t)}(b,v) = \Delta_v - \delta^t_\phi$, and
\[
    \tau_U^t \;:=\; \Delta^t_b \;=\; \tau_U - \delta^t_\phi,
    \qquad
    \tau_L^t \;:=\; \Delta^t_a \;=\; \tau_L - \delta^t_\phi.
\]
\end{definition}
\noindent Both thresholds shift by the same amount, so $\tau_U^t - \tau_L^t = \tau_U - \tau_L$ for all $t$: the dynamic transition interval $(\tau_L^t, \tau_U^t)$ is a rigid translation of the static one, with the same width.  A large positive (resp.\ negative) $\delta^t_\phi$ slides the entire interval down (resp.\ up) the real line.  Because the width is preserved, the drift mechanism that pushes $\Delta Q$ out of the interval (Lemma~\ref{lem:drift}) is unaffected by $\phi$ --- only the \emph{location} of the interval moves.
\begin{theorem}[Dynamic stability of the outgoing message]\label{thm:dynstability}
Fix iteration $t \geq 2$ and let $\tau_L^t, \tau_U^t$ be as in Definition~\ref{def:dynthresh}.
\begin{enumerate}
    \item \textbf{(Upper regime at step $t$.)} If $\Delta Q^{t-2} \geq \tau_U^t$, then
    \[
        \Delta R^t \;=\; C'(a,b) - C'(a,a),
    \]
    independently of $Q^{t-2}_{X_1 \to F_{12}}$ and of $\phi^t$.
    \item \textbf{(Lower regime at step $t$.)} If $\Delta Q^{t-2} \leq \tau_L^t$, then
    \[
        \Delta R^t \;=\; C'(b,b) - C'(b,a),
    \]
    independently of $Q^{t-2}_{X_1 \to F_{12}}$ and of $\phi^t$.
\end{enumerate}
In both cases the value of $\Delta R^t$ is the same constant as in the static Theorem~\ref{thm:stability}; only the gating thresholds depend on $t$.
\end{theorem}
\begin{proof}
Apply Theorem~\ref{thm:stability} to $C'^{(t)}$, which has margins $\Delta^t_v$ and thresholds $\tau_L^t, \tau_U^t$ by construction.  In each case the selected assignment of $X_1$ contributes a constant $\phi^t(\cdot)$ to both values of $X_2$, which cancels in the difference: for the upper regime, $C'^{(t)}(a,b) - C'^{(t)}(a,a) = C'(a,b) - C'(a,a)$, and symmetrically for the lower. \qedhere
\end{proof}
Because $\Delta R$ in either regime is unchanged by $\phi$, the backward step that produces the next $\Delta Q$ uses the same $C'$ as in the static case, and the round-trip constant $c_U$ computed in Section~\ref{sec:absorbing} still applies.
\subsubsection{Generalized Convergence under Bounded Perturbation}\label{subsec:gencvg}
The thresholds $\tau_v^t = \tau_v - \delta^t_\phi$ shift with $\phi$, so the static absorption arguments of Sections~\ref{sec:absorbing}--\ref{sec:convergence} no longer apply directly.  We now bound how much $\phi$ can move them without ejecting the message that has settled inside.
Recall the upper round-trip constant $c_U$ of Section~\ref{sec:absorbing}: once $\Delta Q$ enters the upper regime, it settles to $c_U$, which sits \emph{above} the threshold $\tau_U$.  The gap $c_U - \tau_U$ is how much room the settled message has to spare.  A dynamic unary constraint shifts the upper threshold to $\tau_U - \delta_\phi$ (Definition~\ref{def:dynthresh}); the message stays absorbed as long as $\delta_\phi$ stays within this gap.
By the case analysis in the proof of Theorem~\ref{thm:convergence}, the gap satisfies
\[
    c_U - \tau_U \;\geq\; \min\!\bigl(2|d|,\; -(\tau_U+\tau_L)\bigr),
\]
with equality when $\tau_U \geq -|d|$ (strict only when $\tau_U < -|d|$, where the gap equals $-2\tau_U$).  We therefore require
\[
    B \;<\; \min\!\bigl(2|d|,\; -(\tau_U + \tau_L)\bigr),
\]
which presumes $\tau_U + \tau_L < 0$.  Under this bound, every shift stays strictly within the margin; Theorem~\ref{thm:dyncvg} below confirms that absorption survives.
Note that the bound on $B$ depends on the regime margin ($-(\tau_U + \tau_L)$) and the drift ($2|d|$), but not on the interval width $\tau_U - \tau_L$---the width sets how fast static convergence happens, not how much perturbation the regime can absorb.
\begin{lemma}[Dynamic drift in the transition interval]\label{lem:dyndrift}
Assume $d > 0$ and $|\delta^t_\phi| \leq B$ for all $t$, with $B < \min(2|d|,\,-(\tau_U+\tau_L))$.  For any iteration $t$ with $\Delta Q^t \in (\tau_L^{t+2}, \tau_U^{t+2})$, after one round-trip either
\[
    \Delta Q^{t+4} \;\geq\; \tau_U^{t+6} \qquad\text{(upper regime entered),}
\]
or $\Delta Q^{t+4}$ again lies in a transition interval and the round-trip increment satisfies
\[
    \Delta Q^{t+4} - \Delta Q^t \;\geq\; 2|d| - B \;>\; 0 .
\]
In the split case the increment equals $2|d| + \delta^{t+2}_\phi \in [\,2|d|-B,\,2|d|+B\,]$, so the static rate $2|d|$ is replaced by a bounded \emph{rate window} of width $2B$ centred at $2|d|$; the boundary cases can only increase $\Delta Q$ faster.
\end{lemma}
\begin{proof}
By Theorem~\ref{thm:dynstability}, the forward step at iteration $t+2$ uses the effective cost $C'^{(t+2)}$, with diagonal difference $d + \delta^{t+2}_\phi$.  Since $\Delta Q^t$ is in the moving interval, the factor is in the split regime, so
\[
    \Delta R^{t+2} \;=\; (d+\delta^{t+2}_\phi) + \Delta Q^t \;=:\; r .
\]
The backward step uses the bare $C'$, with the same three regions on $r$ as in the proof of Lemma~\ref{lem:drift}.
\smallskip\noindent\textbf{Both assignments prefer $X_2=a$} ($r \geq -d-\tau_L$): $\Delta Q^{t+4} = -\tau_L$.  Since $\delta^{t+6}_\phi \geq -B > \tau_U+\tau_L$, we have $-\tau_L \geq \tau_U^{t+6}$, so the upper regime is entered.
\smallskip\noindent\textbf{Both prefer $X_2=b$} ($r \leq -d-\tau_U$): $\Delta Q^{t+4} = -\tau_U$.  The gating condition gives $\Delta Q^t \leq -2|d| - \tau_U - \delta^{t+2}_\phi$, so the increment $-\tau_U - \Delta Q^t \geq 2|d| + \delta^{t+2}_\phi \geq 2|d|-B$.  If $-\tau_U \geq \tau_U^{t+6}$ the upper regime is entered; otherwise the iterate remains in the transition interval with the stated increment.
\smallskip\noindent\textbf{Split} (otherwise): $\Delta Q^{t+4} = 2d+\delta^{t+2}_\phi+\Delta Q^t$, an increment of $2|d|+\delta^{t+2}_\phi \in [2|d|-B,\,2|d|+B]$, strictly positive since $B<2|d|$.
\smallskip\noindent At least one of the first two cases is vacuous (both at the boundary $-(\tau_U+\tau_L) = 2|d|$); the increment bound holds in every nonvacuous case. \qedhere
\end{proof}
\begin{theorem}[Generalized convergence under bounded perturbation]\label{thm:dyncvg}
Assume $d > 0$ (the WLOG orientation of Section~\ref{sec:convergence}; the case $d<0$ is symmetric, with the lower regime and constant $c_L$ in place of the upper regime and $c_U$), and let $|\delta^t_\phi| \leq B$ for all $t$, with $B < \min(2|d|,\,-(\tau_U+\tau_L))$.  Starting from any $\Delta Q^0 \in \mathbb{R}$, the trajectory reaches the upper regime in at most
\[
    \left\lceil \frac{\tau_U - \tau_L + 2B}{2|d| - B} \right\rceil
\]
round-trips and remains there for every subsequent iteration; the absorbing value is the static round-trip constant $c_U$ of Section~\ref{sec:absorbing}, given in closed form by the case analysis in the proof of Theorem~\ref{thm:convergence}.
\end{theorem}
\begin{proof}
We treat $d>0$; the case $d<0$ is symmetric.
\textbf{Ascent.}  By Lemma~\ref{lem:dyndrift}, every round-trip whose initial iterate lies in a transition interval either enters the upper regime or strictly increases $\Delta Q$ by at least $2|d|-B>0$, with no decrease.  Hence along round-trips $\Delta Q$ is monotonically increasing until the upper regime is reached.
\textbf{Reaching the upper regime.}  Each interval $(\tau_L^j,\tau_U^j)$ is the static interval translated by $-\delta^j_\phi$ with $|\delta^j_\phi|\leq B$, so its endpoints stay within $[\tau_L-B,\,\tau_U+B]$.  An in-interval iterate is therefore above $\tau_L-B$, and the upper regime is guaranteed once $\Delta Q$ rises above $\tau_U+B$.  Since each in-interval round-trip raises $\Delta Q$ by at least $2|d|-B$, the upper regime is reached after at most
\[
    \left\lceil \frac{(\tau_U+B)-(\tau_L-B)}{2|d|-B} \right\rceil
    = \left\lceil \frac{\tau_U-\tau_L+2B}{2|d|-B} \right\rceil
\]
round-trips.  (If $\Delta Q^0 \leq \tau_L^1$, Lemma~\ref{lem:upper-nonabsorb} gives $\Delta Q^4 = c_L > \tau_L$ after one round-trip; if this is not already in the upper regime, the in-interval bound above applies; cf.\ the proof of Theorem~\ref{thm:convergence}.)
\textbf{Persistence.}  Suppose $\Delta Q^j \geq \tau_U^{j+2}$, so step $j+2$ is in the upper regime and $\Delta R^{j+2} = C'(a,b)-C'(a,a)$ by Theorem~\ref{thm:dynstability}.  The backward update uses the bare $C'$, so it returns the static round-trip constant $c_U$ from Section~\ref{sec:absorbing}, independent of $\phi$.  The next regime test $\Delta Q^{j+4} \geq \tau_U^{j+6}$ reads
\[
    c_U \;\geq\; \tau_U - \delta^{j+6}_\phi
    \quad\Longleftrightarrow\quad
    \delta^{j+6}_\phi \;\geq\; \tau_U - c_U .
\]
By the case analysis in the proof of Theorem~\ref{thm:convergence}, $c_U - \tau_U \geq \min(2|d|,\,-(\tau_U+\tau_L))$, so the condition is implied by $-\delta^{j+6}_\phi \leq B < \min(2|d|,\,-(\tau_U+\tau_L))$.  By induction the upper regime is never left, and $\Delta Q^j = c_U$ thereafter. \qedhere
\end{proof}
The bound interpolates between the static and adversarial limits: as $B \to 0$ it recovers the static $\lceil(\tau_U - \tau_L)/(2|d|)\rceil$ count, and as $B \to \min(2|d|, -(\tau_U+\tau_L))$ both penalties --- the $+2B$ in the numerator (moving thresholds) and the $-B$ in the denominator (partial drift cancellation) --- blow it up.
\begin{remark}[Tightness of the perturbation bound]\label{rem:tight}
The condition $B < \min(2|d|,-(\tau_U+\tau_L))$ is sharp: neither term can be dropped against an arbitrary adversary obeying $|\delta^t_\phi| \leq B$.  If $B \geq 2|d|$, an adversary with $\delta^t_\phi = -2|d|$ on each forward step cancels the drift entirely; if $B > -(\tau_U+\tau_L)$, the moving threshold $\tau_U - \delta_\phi$ can fall below $-\tau_L$ and eject $\Delta Q$ from the upper regime.
\end{remark}
% ------------------------------------------------------------------
\subsection{Toward Damped and Asymmetric Splits}\label{sec:extensions}
% ------------------------------------------------------------------
% NEW in revision. AUTHORS: please verify both propositions carefully
% (the algebra is a direct substitution into Thms 4.5/4.6, but confirm),
% and add them to the Lean development before submission.
Two gaps separate the analysis above from the algorithms evaluated in Section~\ref{Sec:experiments}: DMS-SCFG uses damping, and practical splits need not be symmetric (DABP uses a $0.95/0.05$ split). We take a first step on both fronts.

\paragraph{Damping.}
Damping replaces each sent message with a convex combination of the previous message on the same directed edge and the newly computed one (Section~\ref{Sec:background}). Within a locked regime the newly computed difference is a constant of $C'$ (Theorem~\ref{thm:stability}), so damping acts as exponential smoothing toward that same constant:
\begin{proposition}[Geometric settling under damping]\label{prop:damped}
Fix $\lambda \in [0,1)$ and let messages be damped as in Section~\ref{Sec:background}. Suppose that for all iterations $k \geq t$ the incoming difference satisfies $\Delta Q^{k-2} \geq \tau_U$. Then the damped outgoing difference satisfies
\[
    \Delta R^{k} - K \;=\; \lambda\,\bigl(\Delta R^{k-1} - K\bigr),
    \qquad K := C'(a,b) - C'(a,a),
\]
so $\Delta R^k \to K$ geometrically at rate $\lambda$; the regime constants of Theorem~\ref{thm:stability} are unchanged by damping.
\end{proposition}
\begin{proof}
By Theorem~\ref{thm:stability}, the undamped computation yields the constant difference $\widehat{\Delta R}^{\,k} = K$ whenever $\Delta Q^{k-2} \geq \tau_U$. The damped update on the directed edge $F_{12} \to X_2$ gives $\Delta R^{k} = \lambda\, \Delta R^{k-1} + (1-\lambda)\, \widehat{\Delta R}^{\,k}$; subtracting $K$ from both sides yields the claim.
\end{proof}
Proposition~\ref{prop:damped} is conditional: it assumes the trajectory remains in the upper regime. Closing the induction --- showing that once the damped trajectory enters the regime with sufficient margin it never leaves --- mirrors the persistence argument of Theorem~\ref{thm:dyncvg}, with the geometrically decaying damping residual playing the role of the bounded perturbation $\delta_\phi$. We leave the full damped analogue of Theorem~\ref{thm:convergence} for future work.

\paragraph{Asymmetric splits.}
Let the two copies carry arbitrary tables $C'$ and $C''$ with $C' + C'' = C$, so that the forward update~\eqref{eq:minsumR} uses $C'$ and the backward update~\eqref{eq:minsumQ} uses $C''$. The regime lock of Theorem~\ref{thm:stability} depends only on the forward table, and the round-trip cancellation goes through unchanged:
\begin{proposition}[Asymmetric round-trip constants]\label{prop:asym}
Let $\tau_L, \tau_U$ be the thresholds of Definition~\ref{def:thresholds} computed from $C'$. If $\Delta Q^{t-2} \geq \tau_U$, then $\Delta R^t = C'(a,b) - C'(a,a)$ and
\begin{multline*}
    \Delta Q^{t+2} = \min\!\bigl(C''(b,a)+C'(a,a),\; C''(b,b)+C'(a,b)\bigr) \\
    - \min\!\bigl(C''(a,a)+C'(a,a),\; C''(a,b)+C'(a,b)\bigr) =: c_U(C',C''),
\end{multline*}
independently of $Q^{t-2}$; if additionally $c_U(C',C'') \geq \tau_U$, the upper regime is absorbing for this direction of the cycle. Symmetrically, for $\Delta Q^{t-2} \leq \tau_L$,
\begin{multline*}
    c_L(C',C'') := \min\!\bigl(C''(b,a)+C'(b,a),\; C''(b,b)+C'(b,b)\bigr) \\
    - \min\!\bigl(C''(a,a)+C'(b,a),\; C''(a,b)+C'(b,b)\bigr).
\end{multline*}
Setting $C'' = C'$ recovers~\eqref{eq:cL} and~\eqref{eq:cU}.
\end{proposition}
\begin{proof}
Identical to the proof of Theorem~\ref{thm:absorbing}: in the upper regime $(R^t)_v = C'(a,v) + (Q^{t-2})_a$ by Theorem~\ref{thm:stability}; substituting into the backward update with $C''$ in place of $C'$, the constant shift $(Q^{t-2})_a$ cancels in the difference, and the same induction gives absorption when the constant lies inside the regime. The mirrored flow of the cycle (starting at $X_2$) satisfies the symmetric statement with $C'$ and $C''$ exchanged.
\end{proof}
\begin{remark}[Split-ratio design]\label{rem:ratio}
For a $\rho/(1{-}\rho)$ split, $C' = \rho\, C$ and $C'' = (1{-}\rho)\, C$, and Proposition~\ref{prop:asym} makes the absorption condition $c_U(C',C'') \geq \tau_U(C')$ an explicit function of $\rho$ and the entries of $C$. This turns the split ratio --- treated as a fixed design constant by \citet{CohenGZ20} and \citet{DengKL022} --- into a quantity that can, in principle, be optimized per function-node, and it offers a candidate explanation for the ratio effects observed in the DABP ablation of Section~\ref{Sec:experiments}. Extending the unconditional convergence theorem (Theorem~\ref{thm:convergence}) to the asymmetric case requires redoing the drift case analysis with two tables and is left for future work.
\end{remark}

% ========== END sec4_effect_of_splitting.tex ==========

% ========== BEGIN sec5_splitting_no_damping.tex ==========
\section{Splitting with No Damping}
\label{Sec:MS-split}
The experiments presented in~\cite{CohenGZ20} and in Section~\ref{Sec:experiments} demonstrate that when standard Min-sum is performed on an SCFG, the algorithm oscillates between low quality solutions and, in many cases, between only two such solutions.
\begin{figure}
\centering
\includegraphics[width=8cm]{two_cycle_examp.pdf}
\vspace{-10pt}
\caption{A symmetric SCFG in which two binary function-nodes were split. The tables show the original costs; each split copy holds half of every entry.}
\vspace{-5pt}
\label{Fig:scfg}
\end{figure}
To investigate this behavior, we examined an example of a lemniscate factor graph, which was presented in~\cite{CohenGZ20}. This factor graph, depicted in Figure~\ref{Fig:scfg}, includes two cycles that were generated from the symmetric split of two function-nodes. As stated in~\cite{CohenGZ20}, the MS algorithm does not converge on this factor-graph. In fact, it alternately visits two solutions, one with cost 300 and the other with cost 1200. However, when damping is used, the algorithm converges to the optimal solution with cost 130.
For the case of a single cycle factor graph,~\citet{ForneyKMT2001} prove that the oscillation between multiple solutions is the result of a minimal inconsistent route that the algorithm follows. Thus, one may think that the oscillation in the case of a lemniscate is also caused by an inconsistent minimal route. In fact, after analyzing the case depicted in Figure~\ref{Fig:scfg}, we arrive at the following observation:
\begin{observation}[Unsynchronized period-2 oscillation]\label{obs:lemniscate}
On the lemniscate SCFG of Figure~\ref{Fig:scfg}, MS converges to a period-2 oscillation between two \emph{solutions}: the optimal solution, whose cost is 130, and a disjoint alternative solution with cost 132. The variable-nodes visit the two solutions out of phase --- in iterations in which $X_1$ and $X_3$ take the value assignments of the optimal solution, $X_2$ takes the value assignment of the alternative solution, and vice versa --- so the decoded joint assignments have costs 300 and 1200.
\end{observation}
\noindent This explains why neither of the two good solutions appears in the reported results: the nodes are not synchronized in their oscillation between them. (A BCT view of the same phenomenon --- the alternating structure of the bottom layers of the middle variable's backtrack cost tree --- is consistent with this observation, but the general account is given in Section~\ref{sec:tworoutes}: once the split function-nodes commit, the decoded assignments follow a synchronous local selection rule whose only limit behaviors are convergence and an alternation between exactly two solutions, and the out-of-phase structure of Observation~\ref{obs:lemniscate} is precisely the bipartite rephasing of Corollary~\ref{cor:bipartite}, under which the two interleaved snapshots of costs 300 and 1200 carry two locally optimal solutions of costs 130 and 132.)
Building on Observation~\ref{obs:lemniscate}, we propose two versions of the algorithm that exploit the two-branch structure of the oscillation. Given a factor graph, we generate its SCFG and perform 200 MS iterations. We take the assignments selected by the algorithm at iterations 198 and 199 --- the two branches of the (empirically period-2) oscillation --- and generate a new problem that includes in the domain of each variable only these two values. We then use either MGM~\cite{MaheswaranTBPV04} or SyncBB~\cite{HirayamaY97} to select a solution of this reduced problem. As the experiments in Section~\ref{Sec:experiments} show, the improvement in the quality of the selected solutions, relative to the raw MS-SCFG output, is dramatic and statistically significant on every benchmark (Table~\ref{tab:merge}); as expected, however, these undamped variants do not match the split-damped algorithms on the structured benchmarks (cf.\ Table~\ref{tab:final_costs}).

% ========== END sec5_splitting_no_damping.tex ==========

% ========== BEGIN sec5b_two_solutions.tex ==========
% ======================================================================
% Main-text version of the two-solution theory. Full proofs of
% Lemmas ~\ref{lem:clones}, \ref{lem:decode}, \ref{lem:ema},
% Theorem~\ref{thm:lambdastar}, Proposition~\ref{prop:kn}, and
% Corollary~\ref{cor:bipartite}, together with the measurement
% protocol, census, and audits, are in Appendix~\ref{app:tworoutes}.
% ======================================================================
\section{The Two-Solution Structure of Min-sum on SCFGs}\label{sec:tworoutes}

We now explain the behavior documented in Section~\ref{Sec:MS-split} in
three conditional steps: (i) a split function-node whose incoming
difference is in a regime transmits a \emph{decision} --- one row of its
cost table --- rather than a magnitude, and when this holds for every
message involved (\emph{full commitment}), the selected assignments evolve
by a simple synchronous rule; (ii) \emph{any} assignment sequence
following that rule is eventually periodic with period one or two, and the
limiting pair is a local minimum of a pair cost we call the alternation
cost; (iii) damping destroys a fixed committed period-2 pattern exactly
above an explicit threshold, and the threshold approaches $1$ as the
degree grows, so no single damping value suffices universally
(Section~\ref{sec:whydamping}).  Throughout, $x$ and $y$ denote complete
assignments, $x_i \in D_i$ is the value $x$ assigns to $X_i$, and
$\phi_i : D_i \to \mathbb{R}$ denotes the small unary tie-breaking
preferences used in the experiments (the dynamic unary constraint
$\phi^t$ of Section~\ref{sec:dynamic} is the time-varying case).  Scalar
statements (regimes, transition intervals) are for binary domains; claims
about larger domains are labeled empirical.

\begin{remark}[Schedule and orientation]\label{rem:conv}
We index messages by the two-phase schedule of
Section~\ref{secmaxsum}: in iteration $t$ all variable-nodes send
$Q^{t}$ (computed from $R^{t-1}$), then all function-nodes send $R^{t}$
(computed from $Q^{t}$); value selection at iteration $t$ uses $R^{t}$.
For cost terms, $C_{ij}$ denotes the \emph{shared} cost term of the
constraint between $X_i$ and $X_j$, read with $X_i$'s value first:
$C_{ij}(v, w) = C_{ji}(w, v)$ by definition.  No numerical symmetry of
the table is assumed.
\end{remark}

\subsection{From Regimes to Synchronous Local Selection}\label{sec:rule}

By Theorem~\ref{thm:stability}, outside the transition interval the
minimum in the $R$-update is attained at the same value of the sending
variable for both values of the receiving variable, so the $R$-message
equals one row of the cost table plus a constant.  For general domains we
say $Q^{t}_{X_i \to F}$ has a \emph{stable active minimizer} $u^{*}$ if
$u^{*} = \operatorname{argmin}_{u}[\,C'(u,v) + (Q^{t}_{X_i\to F})[u]\,]$
is the same for every $v$; then $R^{t}_{F\to X_j}(\cdot) =
C'(u^{*},\cdot) + (Q^{t})[u^{*}]$: \emph{the function-node forwards row
$u^{*}$}, up to an additive constant that shifts all entries equally
(this constant keeps propagating, which matters for damping).

\begin{lemma}[Clone synchronization]\label{lem:clones}
Under the symmetric split with zero-initialized messages, the two copies
$F'_{ij}, F''_{ij}$ send identical messages at every iteration, and so
does $X_i$ toward the two copies.
\end{lemma}

Lemma~\ref{lem:clones} exposes what splitting changes: the message $X_i$
sends toward $F'_{ij}$ \emph{contains the sibling copy's previous
$R$-message}.  Ordinary Min-sum is non-backtracking, whereas on the SCFG
every edge carries a length-two feedback cycle through the sibling copy.
A round trip through this cycle applies the same cost term twice, so it
reinforces its own direction whether the constraint prefers agreement or
disagreement, and its drift per round trip is the constant $d$ of
Section~\ref{sec:convergence}; Lemma~\ref{lem:drift} and
Theorem~\ref{thm:dyncvg} state that the cycle expels $\Delta Q$ from the
transition interval and holds it outside when the external swing obeys
the stated bound.  On benchmark graphs commitment is partial and
family-dependent: measured mean fractions of messages with a stable
active minimizer at iteration $9$ / in the tail are $0.73$/$0.82$
(random sparse, $|D_i|{=}10$, split) versus $0.12$/$0.10$ (unsplit),
$0.74$/$0.92$ versus $0.36$/$0.35$ (random dense), $0.97$/$0.97$ versus
$0.87$/$0.85$ (random binary), and $0.03$/$0.36$ versus $0.00$/$0.00$
(graph coloring, which rarely commits in this strict sense although its
messages still carry the sender's selection; Appendix~\ref{app:tworoutes}).
Commitment can also hold on one parity of an alternation and fail on the
other, and for $|D_i| > 2$ the commitment region is not an interval.

\begin{lemma}[Decoded assignment under full commitment]\label{lem:decode}
If at iteration $t$, for every neighbor $X_j$ of $X_i$, the message
$Q^{t}_{X_j \to F'_{ij}}$ has a stable active minimizer $u^{*t}_j$, then
\[
  \hat{x}^{t}_i \;=\; \operatorname{argmin}_{v \in D_i}
  \Big[\phi_i(v) + \sum_{j : C_{ij}\in\mathcal{C}}
  C_{ij}\big(v,\, u^{*t}_j\big)\Big].
\]
\end{lemma}

\begin{lemma}[Sibling evaluation]\label{lem:sibling}
Assume the hypothesis of Lemma~\ref{lem:decode} at iteration $t$, and
additionally that $Q^{t+1}_{X_i \to F'_{ij}}$ has a stable active
minimizer $u^{*,t+1}_i$, with unique argmins.  Then
$u^{*,t+1}_i = \hat{x}^{t}_i$.
\end{lemma}
\begin{proof}
By Lemma~\ref{lem:clones} and the hypothesis at $t$,
$(Q^{t+1}_{X_i \to F'_{ij}})[u] = \phi_i(u) + \sum_{k \neq j}
C_{ik}(u, u^{*t}_k) + \tfrac12 C_{ij}(u, u^{*t}_j) + \text{const}$, the
last term being the sibling's forwarded half-row.  A stable active
minimizer minimizes $\tfrac12 C_{ij}(u,v) + (Q^{t+1})[u]$ for
\emph{every} $v$; evaluating at $v = u^{*t}_j$ makes the two half-costs
of $C_{ij}$ recombine, so the minimized expression becomes the full field
$\phi_i(u) + \sum_{k} C_{ik}(u, u^{*t}_k) + \text{const}$, whose unique
argmin is $\hat{x}^{t}_i$ by Lemma~\ref{lem:decode}. \qedhere
\end{proof}

\begin{corollary}[The selection rule]\label{cor:rule}
If the hypotheses of Lemmas~\ref{lem:decode}--\ref{lem:sibling} hold for
every variable and constraint at every $t \ge t_0$ (\emph{full
commitment}), then for all $t \ge t_0$, simultaneously for all $i$,
\[
  \hat{x}^{t+1}_i \;=\; \operatorname{argmin}_{v\in D_i}
  \Big[\phi_i(v) + \sum_{j : C_{ij}\in\mathcal{C}}
  C_{ij}\big(v, \hat{x}^{t}_j\big)\Big]
\]
--- \emph{synchronous local selection}.
\end{corollary}

Because full commitment is partial in practice, we state the attractor
theorem under the weaker hypothesis that the selected assignments satisfy
this recursion --- implied by full commitment, and checkable directly on
a run.  On an audited subset of $32$ period-2 runs across four benchmark
families, the recursion held on every check performed: the two
alternating solutions mapped to one another under it exactly (both
directions, on the limit pair of every instance), and it predicted
$\hat{x}^{t+1}$ from $\hat{x}^{t}$ correctly on all sampled transitions
(every seventh transition of the recorded tails; protocol in
Appendix~\ref{app:tworoutes}).

\subsection{The Two-Solution Theorem}\label{sec:twosol}

\begin{definition}[Alternation cost]\label{def:altcost}
For complete assignments $(x, y)$ let
\[
  \mathrm{cost}_2(x,y) = \!\!\sum_{C_{ij}\in\mathcal{C}}\!\!
  \big[C_{ij}(x_i, y_j) + C_{ij}(y_i, x_j)\big]
  + \sum_i \big[\phi_i(x_i) + \phi_i(y_i)\big].
\]
Then $\mathrm{cost}_2(x,y) = \mathrm{cost}_2(y,x)$ --- using only that
the two endpoints share one cost term --- and $\mathrm{cost}_2(x,x) =
2\,\mathrm{cost}(x)$.
\end{definition}

\begin{theorem}[Two solutions]\label{thm:tworoutes}
If for all $t \ge t_0$ the selected assignments satisfy the synchronous
local selection recursion with unique argmins, then
$\mathrm{cost}_2(\hat{x}^{t}, \hat{x}^{t+1})$ is non-increasing and
strictly decreasing unless $\hat{x}^{t+2} = \hat{x}^{t}$; hence the
selected-assignment sequence is eventually periodic with period $1$ or
$2$.
\end{theorem}
\begin{proof}
Grouping $\mathrm{cost}_2(\hat{x}^{t+1}, z)$ around $z$, each $X_i$
contributes $\phi_i(z_i) + \sum_{j} C_{ij}(z_i, \hat{x}^{t+1}_j)$ plus
terms independent of $z$; this is separable and coordinatewise minimized
by $z = \hat{x}^{t+2}$.  Hence
$\mathrm{cost}_2(\hat{x}^{t+1}, \hat{x}^{t+2}) \le
\mathrm{cost}_2(\hat{x}^{t+1}, \hat{x}^{t}) =
\mathrm{cost}_2(\hat{x}^{t}, \hat{x}^{t+1})$ by swap symmetry.  With
unique argmins, equality forces $\hat{x}^{t+2} = \hat{x}^{t}$;
$\mathrm{cost}_2$ takes finitely many values, so strict decreases
terminate. \qedhere
\end{proof}

\begin{corollary}\label{cor:localmin}
The limit pair $(x, y)$ is a local minimum of $\mathrm{cost}_2$ under
single-variable changes in either coordinate.  Period $1$ occurs exactly
when the pair is diagonal, and a diagonal limit is a local minimum of the
original problem.
\end{corollary}

Parallel-update period-at-most-two results go back to
\citet{GolesOlivos80} (threshold networks, symmetric weights) and
\citet{PoljakSura83} (weighted-plurality opinion dynamics --- pairwise
terms proportional to equality indicators); a two-cycle result for
simultaneous best response in two-player random potential games appears
in \citet{AshkenaziGolan25}.  Theorem~\ref{thm:tworoutes} covers
arbitrary \emph{shared} pairwise cost terms with unaries, hence the short
proof.  Within its hypothesis it explains why the observed period is two
rather than three, five, or aperiodic.  Empirically (census in
Appendix~\ref{app:tworoutes}; periods detected on the final $150$ of
$700$ iterations, up to $64$): of $276$ split runs, $274$ had a detected
periodic tail of period $1$ or $2$ ($191$ period-2, $83$ period-1) and
two had longer detected periods ($16$, $42$); restricted to the
$n \ge 40$ cells of the domain-10 and coloring families the period-2
count is $92$ of $92$; unsplit runs on the random domain-10 families had
\emph{no} detected period at all in $85$--$100\%$ of instances from
$n{=}20$ up.

Off-diagonal minima of $\mathrm{cost}_2$ are where alternation pays: for
three mutually ``not-equal'' constrained variables, every single
assignment must violate one constraint, but the pair
$x = (\alpha,\alpha,\alpha)$, $y = (\beta,\beta,\beta)$ satisfies all six
cross terms.  On the audited subset, the limit pair was a local minimum
of $\mathrm{cost}_2$ with zero violating moves in $32/32$ runs.

\begin{corollary}[Bipartite rephasing]\label{cor:bipartite}
Let the constraint graph be connected and bipartite with sides $A, B$,
and $(x,y)$ any pair.  Define $x'$ ($=x$ on $A$, $=y$ on $B$) and $y'$
($=y$ on $A$, $=x$ on $B$).  Then
$\mathrm{cost}_2(x, y) = \mathrm{cost}(x') + \mathrm{cost}(y')$, and if
$(x,y)$ is a layer-wise local minimum of $\mathrm{cost}_2$ then $x'$ and
$y'$ are each coordinatewise local minima of the \emph{original} problem.
\end{corollary}

On a bipartite graph, then, the oscillating algorithm is not wandering
between two poor assignments: it holds \emph{two locally optimal
solutions in interleaved form}, and the raw snapshots read them out in
mixed phase.  This is exactly Observation~\ref{obs:lemniscate}
($\{X_1, X_3\}$ and $\{X_2\}$ are the two sides; snapshots cost $300$ and
$1200$ while their rephasings are the local optima of costs $130$ and
$132$), and it was verified on $8/8$ oscillating complete-bipartite
random instances (Appendix~\ref{app:tworoutes}).  On non-bipartite graphs
the only exact symmetry is the global swap $(x,y)\mapsto(y,x)$;
per-variable rephasing changes $\mathrm{cost}_2$ and is heuristic
post-processing (measured gains in Appendix~\ref{app:tworoutes}).  The
two-value menus of the merge heuristics of Section~\ref{Sec:MS-split} are
exactly the coordinates of the limit pair: on bipartite graphs
Corollary~\ref{cor:bipartite} guarantees they contain two coordinatewise
local optima; on general graphs the support is empirical
(Table~\ref{tab:merge}).

\section{Damping and the Two-Solution Oscillation}\label{sec:whydamping}

Damping is applied to the $Q$-messages
(Section~\ref{secmaxsum}).  Because a committed function-node forwards
its row plus a propagating constant, raw messages need not settle even
when all decisions are locked; the objects that settle are the message
\emph{differences}.  All statements below are about differences.

\begin{lemma}[Averaging an alternating computation]\label{lem:ema}
If a damped scalar sequence has computed values alternating between
constants $u$ and $v$, the sent values converge geometrically (rate
$\lambda$) to the alternating pair $\frac{u+\lambda v}{1+\lambda}$,
$\frac{v+\lambda u}{1+\lambda}$: the mean is preserved and the swing is
multiplied by $(1-\lambda)/(1+\lambda)$ --- $1/19$ at $\lambda = 0.9$.
\end{lemma}

\begin{theorem}[Persistence threshold of a committed alternation]\label{thm:lambdastar}
Binary domains.  Consider a trajectory of $Q$-damped Min-sum on the SCFG
that, from some iteration on, reproduces a fixed period-2 pattern of
regimes in which every message difference is in a regime at every step.
For a message whose regime side flips each iteration, the computed
differences alternate between constants $\Delta_1 > \tau_U$ and
$\Delta_2 < \tau_L$ determined by the pattern, and the pattern can
persist only if
\[
  \lambda \;\le\;
  \min\!\left(
    \frac{\Delta_1 - \tau_U}{\tau_U - \Delta_2},\;
    \frac{\tau_L - \Delta_2}{\Delta_1 - \tau_L}
  \right)
\]
for every flipping message; conversely, under strict inequalities the
corresponding period-2 orbit of damped message differences exists and is
locally attracting (contraction $\lambda$ per step).  The bound is
strictly below $1$ whenever $\tau_U > \tau_L$ or the two margins differ.
Damped and undamped updates have identical fixed points.
\end{theorem}

What Theorem~\ref{thm:lambdastar} does \emph{not} say: above the
threshold of a given pattern the trajectory need not converge ---
message differences may enter transition intervals, settle into a
different pattern, or sustain non-committed oscillations.  Indeed, of
$17$ initially oscillating instances in our damping scans, $2$ (both
random dense) were non-monotone in $\lambda$.  We therefore report per
instance the smallest grid $\lambda$ from which all larger tested values
produced a constant \emph{assignment} tail (final $300$ of $2000$
iterations; assignment constancy, cost constancy, and convergence of
message differences are increasingly strong statements and only the
first was tested): at most $0.2$ on random sparse and binary, with the
largest values, $0.4$, on three coloring instances and one random dense
instance (Appendix~\ref{app:tworoutes}).  For $|D_i| > 2$ the theorem
does not apply and the domain-10 results are empirical.

\begin{proposition}[The threshold approaches $1$ with degree]\label{prop:kn}
Let the constraint graph be complete on $n$ agents with binary domains,
every constraint costing $10$ when the two agents take equal values and
$0$ otherwise, symmetric split, and a uniform unary preference gap
$0.01$.  By symmetry all message differences share one value $q_t$,
whose damped dynamics is the scalar recurrence
\[
  q_{t+1} \;=\; \lambda q_t + (1-\lambda)\big[\,0.01 -
  (2n-3)\operatorname{clip}(q_t, -5, 5)\,\big],
\]
and the alternation between the two uniform solutions is a fully
committed period-2 pattern with persistence threshold
\[
  \frac{10(n-2) - 0.01}{10(n-1) + 0.01} \;\xrightarrow[n\to\infty]{}\; 1 .
\]
Hence for every fixed $\lambda < 1$ there is a degree at which the
oscillation persists: no problem-independent damping value exists.  For
$n = 12$ the threshold is $0.9089$; the customary $\lambda = 0.9$ does
not eliminate this oscillation (verified in the implementation from zero
initialization), while at $\lambda = 0.95$ the message differences
converge.
\end{proposition}

\paragraph{Summary, with scope.}
Under full commitment and clone synchronization, split Min-sum's
selections provably follow synchronous local selection
(Corollary~\ref{cor:rule}); any selection-following tail is eventually of
period $1$ or $2$ and parks at a local minimum of $\mathrm{cost}_2$
(Theorem~\ref{thm:tworoutes}), whose bipartite rephasings are ordinary
local optima (Corollary~\ref{cor:bipartite}).  On the benchmarks,
commitment is partial and family-dependent, yet the selection recursion
held on every check of the audited runs, and detected split tails were
overwhelmingly of period $1$ or $2$ --- period $2$ almost exclusively in
the large domain-10 and coloring cells --- while unsplit runs mostly
showed no detectable period.  For damping: alternating computed
differences are attenuated by $(1-\lambda)/(1+\lambda)$
(Lemma~\ref{lem:ema}); each fully committed binary period-2 pattern
persists exactly up to an explicit threshold
(Theorem~\ref{thm:lambdastar}); and the threshold approaches $1$ with
degree (Proposition~\ref{prop:kn}), so the reliability of
$\lambda = 0.9$ on the benchmark families --- measured stable-damping
values at most $0.4$ --- is a property of those families, not a universal
guarantee.  Damping removes the parallel-update alternation that
committed dynamics admits as its principal failure mode; it cannot
manufacture commitment, which is splitting's contribution --- partial,
but decisive on these benchmarks.

% ========== END sec5b_two_solutions.tex ==========

% ========== BEGIN sec6_experiments.tex ==========
\section{Experimental Evaluation}
\label{Sec:experiments}
Our empirical study is designed to test the mechanism suggested by the
theoretical analysis.  Splitting a function-node does not change the objective
function, but it changes the message-passing dynamics by modifying the local
active-minimizer geometry and by introducing short feedback loops between the
copies of each split function-node.  These feedback loops can separate
candidate assignments quickly, but, without damping, they may also generate
oscillations.  Damping is therefore expected to stabilize the oscillatory
component created by splitting.
The experiments are divided into four parts.  First, we compare standard
Min-sum (MS), damped Min-sum (DMS), Min-sum on split constraint factor
graphs (MS-SCFG), and DMS on split constraint factor graphs (DMS-SCFG).
Second, we evaluate delayed splitting, in which DMS is performed on the
original factor graph for a fixed number of iterations before the function-nodes
are split.  Third, we evaluate post-processing of the two assignments produced
by the split-only oscillation.  Fourth, we compare against Deep Attentive
Belief Propagation (DABP) and two ablations: DABP without splitting and DABP
with a symmetric split.
\subsection{Experimental Setup}
We evaluated the algorithms on five binary benchmark families, generating 50
independent instances per family (seeds $0$--$49$).
\emph{Random sparse} and \emph{random dense}: 50 agents, one variable each
with $|D_i| = 10$, Erd\H{o}s--R\'enyi constraint graphs with density $p_1 = 0.1$
(sparse) and $p_1 = 0.6$ (dense), and integer constraint costs drawn uniformly
from $[100, 200)$; disconnected components are force-connected.
\emph{Graph coloring}: 50 agents, 3 colors, $p_1 = 0.1$, with not-equal
constraints costing 10 when neighboring agents select the same color and 0
otherwise.
\emph{Scale-free}: Barab\'asi--Albert graphs following \citet{CohenGZ20}: 7
initial randomly connected agents, each new agent attaching preferentially to 3
existing agents, $n = 50$, $|D_i| = 10$, integer costs uniform in $[100, 200)$.
\emph{Meeting scheduling}: following \citet{CohenGZ20}, 90 agents schedule 20
meetings into 20 time slots; each agent participates in two random meetings;
each constrained meeting pair is assigned a travel time drawn uniformly from
$\{6, \ldots, 10\}$, and the cost counts the overbooked (shared) agents when
the slot difference is below the travel time.
All instances add tiny random unary value preferences (uniform in
$[0, 10^{-2})$) for tie breaking, as in prior Max-sum studies
\cite{FarinelliRPJ08}; their total mass stays below the smallest structural
cost gap, so they cannot change which assignment is optimal.
All damped variants use $\lambda = 0.9$.

All algorithms were executed for 2000 synchronous iterations.  At each
iteration, the assignment selected by the current beliefs was evaluated using
the original cost tables.  Thus, even when the algorithm runs on a split
constraint factor graph, all costs are measured with respect to the original
problem.  We report means over the 50 instances per family; standard
deviations accompany all summary statistics, and all pairwise comparisons are
evaluated with paired t-tests and Wilcoxon signed-rank tests over the common
instances.  Full per-benchmark summary and significance tables are included
with the code.
We focus on final solution cost and on the cost trajectory over time.  This is
different from the anytime emphasis of previous work.  The goal here is not
only to determine whether an algorithm visits a good assignment, but whether
the combination of splitting and damping stabilizes to a high-quality assignment
by the end of the run.
The DABP curves are plotted on the same iteration axis after a wall-clock
normalization.  In the plot legend, a label of the form \texttt{x6.3} means that
one iteration of that DABP variant took 6.3 times as long as one DMS iteration
on the same machine.  All such timing measurements were taken on the same
GPU workstation with an NVIDIA RTX 4090 24GB GPU.  Therefore, a DABP curve
labelled \texttt{x6.3} should be read as a wall-clock-adjusted curve, not as a
plain per-iteration curve.
\subsection{Algorithms}
We compare the following algorithms; the labels in parentheses are the ones
used in the plot legends.
\textbf{MS.}  Standard synchronous Min-sum on the original factor graph.
\textbf{DMS.}  Damped Min-sum on the original factor graph.
\textbf{MS-SCFG} (\texttt{MS s=.5})\textbf{.}  Standard Min-sum on a symmetric
split constraint factor graph.  Each original function-node is replaced by two
function-nodes whose cost tables each contain half of the original cost table.
\textbf{DMS-SCFG} (\texttt{DMS s=.5})\textbf{.}  Damped Min-sum on the
symmetric split constraint factor graph.
\textbf{DMS-SCFG-rand} (\texttt{DMS s=.4-.6})\textbf{.}  Damped Min-sum on a
\emph{random} SCFG: each cost-table entry $c$ is split into $u \cdot c$ and
$(1-u) \cdot c$ with $u \sim U[0.4, 0.6)$, the best-performing split reported
by \citet{CohenGZ20}.
\textbf{DMS-SCFG-at-$K$} (\texttt{DMS @$K$})\textbf{.}  Delayed splitting.  The
algorithm performs DMS on the original factor graph for $K$ iterations and then
converts the graph to an SCFG, redistributing the previously sent messages
across the two copies.  We evaluate $K\in\{50,100,300,500,1000\}$, and for
random dense graphs we also evaluate $K=1500$.
\textbf{MS-SCFG-MGM-200 and MS-SCFG-opt-200} (\texttt{MGM@200},
\texttt{Opt merge@200})\textbf{.}  These versions use the two assignments
selected at iterations 198 and 199 of the MS-SCFG run --- the two branches of
its period-2 oscillation (Section~\ref{Sec:MS-split}).  The MGM version applies
an MGM-style local improvement restricted to the two values suggested for each
variable by these assignments, started from both branches, keeping the better
result.  The optimal version solves the corresponding restricted binary-menu
problem with branch and bound (300-second cap); the search is exact on every
benchmark except random dense, where the reported value is the best merge found
within the cap, warm-started from the MGM result.  A diagnostic variant
(\texttt{MGM inv@200}) flips every disagreement variable of the MGM merge to
the other value of its binary menu.
\textbf{Optimal} (\texttt{Opt n=$k$})\textbf{.}  Branch and bound on the
original problem with a 60-second limit per instance, reported only for the $k$
instances on which the search completed (21 of 50 on graph coloring; the search
does not finish on the domain-10 benchmarks).
\textbf{DABP.}  Deep Attentive Belief Propagation~\cite{DengKL022}.  DABP
learns dynamic damping factors and neighbor weights for composing BP messages.
The original DABP implementation also uses an SCFG representation with a
$0.95/0.05$ split ratio.
\textbf{DABP-NoSplit.}  The same attentive architecture, but run on the original
factor graph without splitting.  This ablation tests whether learned damping
and learned neighbor weights are sufficient without the split representation.
\textbf{DABP-SymSplit.}  The same attentive architecture, but using a symmetric
$0.5/0.5$ split instead of the $0.95/0.05$ split.  This ablation tests whether
the advantage of DABP depends on the highly asymmetric split or only on the
presence of a split representation.
\subsection{Final Solution Cost}
Table~\ref{tab:final_costs} reports the mean final solution cost after 2000
iterations.  Lower values indicate better solutions.
\begin{table*}[t]
\centering
\caption{Mean final solution cost after 2000 iterations (50 instances per
benchmark; lower is better; best in bold).  Standard deviations and full
pairwise significance tables accompany the code.}
\label{tab:final_costs}
\scriptsize
\begin{tabular}{lrrrrrrrr}
\toprule
\textbf{Problem}
& \textbf{MS}
& \textbf{DMS}
& \textbf{MS-SCFG}
& \textbf{DMS-SCFG}
& \textbf{Best delayed split}
& \textbf{DABP}
& \textbf{DABP-NoSplit}
& \textbf{DABP-SymSplit} \\
\midrule
Graph coloring
& $666.19$
& $138.81$
& $668.61$
& $34.42$
& $\mathbf{30.62}\;(K{=}300)$
& $39.24$
& $136.41$
& $35.02$ \\
Meeting scheduling
& $87.77$
& $20.70$
& $81.80$
& $8.10$
& $\mathbf{7.82}\;(K{=}100)$
& $27.22$
& $24.00$
& $30.15$ \\
Random dense
& $108076.51$
& $101904.84$
& $107330.16$
& $101571.10$
& $101335.62\;(K{=}1000)$
& $\mathbf{99223.68}$
& $100309.79$
& $99665.79$ \\
Random sparse
& $17611.47$
& $15749.35$
& $17288.43$
& $15357.87$
& $15330.27\;(K{=}1000)$
& $\mathbf{14519.01}$
& $15108.47$
& $14631.11$ \\
Scale-free
& $19924.49$
& $18229.40$
& $19703.28$
& $17952.74$
& $17900.34\;(K{=}1000)$
& $\mathbf{16628.46}$
& $17088.25$
& $16771.71$ \\
\bottomrule
\end{tabular}
\end{table*}
The results show that splitting alone is not sufficient.  MS-SCFG remains close
to MS on the structured benchmarks and is much worse than the split-damped
variants.  This supports the theoretical observation that splitting creates a
new feedback structure but may still oscillate without damping.
Damping alone improves MS substantially, but it does not fully explain the
success of split-damped message passing.  In graph coloring, DMS reaches a
mean final cost of $138.81$, while DMS-SCFG reaches $34.42$ and delayed
splitting reaches $30.62$.  In meeting scheduling, DMS reaches $20.70$, while
DMS-SCFG reaches $8.10$ and delayed splitting reaches $7.82$.  These large
gaps are consistent with the claim that splitting changes the local message
geometry in a way that damping alone does not.
The DABP ablation further strengthens this conclusion.  On graph coloring,
DABP-NoSplit reaches $136.41$, close to DMS and far from DABP with splitting
($39.24$) and DMS-SCFG ($34.42$).  Thus, learned damping and learned neighbor
weights do not replace splitting on this structured benchmark.  On meeting
scheduling, both DABP and DABP-NoSplit are much worse than DMS-SCFG and the
best delayed split.  On random sparse, random dense, and scale-free graphs,
DABP with splitting is consistently better than DABP-NoSplit, and DABP attains
the best final costs overall.  The advantage of the split-damped Min-sum
variants is therefore concentrated on the structured benchmarks; on the
unstructured families, DABP's learned per-edge parameters win, at the price of
centralized execution and a substantially higher per-iteration cost.  Overall,
the ablation indicates that splitting remains an essential structural component
even when the algorithm learns dynamic damping factors and neighbor weights.

\paragraph{Statistical significance.}
All comparisons are paired per instance ($n = 50$); we cite Wilcoxon
signed-rank p-values on final costs, and the paired t-tests agree unless noted.
Adding splitting to DMS is significant on graph coloring ($p < 10^{-9}$),
meeting scheduling ($p < 10^{-8}$), random sparse ($p < 10^{-5}$), and
scale-free graphs ($p = 0.018$), but not on random dense graphs ($p = 0.47$).
Removing splitting from DABP is significant on graph coloring
($p < 10^{-9}$), random sparse ($p < 10^{-3}$), and scale-free graphs
($p < 0.005$), but not on meeting scheduling ($p = 0.40$); on random dense
the t-test is significant ($p = 0.003$) while the Wilcoxon test is not
($p = 0.80$).  On graph coloring the difference between DMS-SCFG and DABP is
\emph{not} significant ($p = 0.31$); on meeting scheduling DMS-SCFG is
significantly better than DABP ($p < 10^{-6}$), and on the three unstructured
families DABP is significantly better ($p < 10^{-14}$).  Delayed splitting
improves significantly on DMS-SCFG only on meeting scheduling ($K{=}100$,
$p = 0.031$) and random dense ($K{=}1000$, $p = 0.029$); on graph coloring the
evidence is mixed (t-test $p = 0.038$, Wilcoxon $p = 0.14$), and on random
sparse and scale-free graphs the improvement in the mean is not significant.
\subsection{Cost Trajectories on Structured Benchmarks}
Figure~\ref{Fig:structured_cost_curves} shows the cost trajectories for graph
coloring and meeting scheduling.  These two benchmarks are the most important
for the main paper because they exhibit the clearest separation between the
mechanisms.  In both cases, MS and MS-SCFG remain high or unstable, while
DMS-SCFG and delayed splitting stabilize at much lower costs.  The DABP-NoSplit
curve is also informative: it tests whether neural attention alone can replace
splitting.  The answer is negative on graph coloring, where removing splitting
keeps DABP close to the DMS region rather than the split-based region.
\begin{figure*}[t]
\centering
\includegraphics[width=.49\textwidth]{plots/graph_coloring_cost.pdf}
\includegraphics[width=.49\textwidth]{plots/meeting_scheduling_cost.pdf}
\caption{Mean solution cost per iteration on the structured benchmarks.
Left: graph coloring.  Right: meeting scheduling.  The DABP labels include a
wall-clock factor: for example, \texttt{DABP x6.3} means that this variant is
6.3 times slower per plotted iteration than DMS on the same RTX 4090 24GB
machine.}
\label{Fig:structured_cost_curves}
\end{figure*}
Because several successful algorithms become close after the transient phase,
Figure~\ref{Fig:structured_cost_curves_zoom} shows a zoomed view of the
lower-cost region.  These plots should be read as convergence plots rather than
anytime plots.  They show whether the decoded solution cost remains low near
the end of the run.
\begin{figure*}[t]
\centering
\includegraphics[width=.49\textwidth]{plots/graph_coloring_cost_zoom.pdf}
\includegraphics[width=.49\textwidth]{plots/meeting_scheduling_cost_zoom.pdf}
\caption{Zoomed view of the lower-cost region on graph coloring and meeting
scheduling.  The zoom highlights the separation between DMS-SCFG, delayed
splitting, DABP, DABP-NoSplit, and DABP-SymSplit after the early transient.}
\label{Fig:structured_cost_curves_zoom}
\end{figure*}
\subsection{The Symmetric-Split DABP Ablation}
The DABP-SymSplit line has a specific interpretation.  It keeps the DABP
architecture and online learning procedure fixed, but replaces the $0.95/0.05$
split with a symmetric $0.5/0.5$ split.  Thus it isolates the effect of the split
ratio.
The $0.95/0.05$ split is significantly better than the symmetric split on the
three unstructured families (random sparse, random dense, and scale-free;
Wilcoxon $p < 10^{-8}$ in each case), while on the two structured benchmarks
the two split ratios are statistically indistinguishable (graph coloring
$p = 0.44$, where the symmetric split has the better mean, $35.02$ vs.\
$39.24$; meeting scheduling $p = 0.23$).
The original $0.95/0.05$ split gives one copy almost the entire original cost
table and gives the second copy only a small fraction.  Early in the run, the
dominant copy is therefore likely to select active minimizers similar to those
selected by the original unsplit function-node.  At the same time, the small
copy already creates an additional feedback channel.  Later in the run, this
weak channel can become more important through accumulated messages and
through the learned DABP neighbor weights.  This provides a plausible mechanism
for why the $0.95/0.05$ split yields better final costs on the unstructured
families: it preserves the original factor behavior early, while leaving a
learned secondary route that can be amplified later.  Proposition~\ref{prop:asym}
gives the corresponding formal handle: the split ratio moves the round-trip
constants and thresholds, and hence the absorption condition.
The symmetric split behaves differently.  Both copies immediately operate at
equal scale.  This makes the feedback more balanced from the first iterations
and removes the strong-copy/weak-copy asymmetry.  The symmetric split is
often faster in wall-clock factor on the structured and sparse benchmarks: the
measured per-iteration ratios give graph coloring \texttt{DABP x6.3} versus
\texttt{DABP sym-split x6.1}, meeting scheduling \texttt{x8.2} versus
\texttt{x7.3}, and random sparse \texttt{x6.9} versus \texttt{x4.6}.  The
effect is not universal: on random dense and scale-free graphs the symmetric
split is not faster by this timing factor.  Since the two split DABP variants
have the same number of factor-node copies, this runtime difference should not
be explained by graph size alone.  The likely explanation is that the split
ratio changes the learned message dynamics and the early stopping / convergence
behavior, rather than the raw number of messages.
This ablation suggests a useful future direction.  We should inspect the learned
DABP neighbor weights and $R$-message magnitudes on the two split copies.  If
the interpretation above is correct, then the $0.95/0.05$ split should show a
weak copy that is initially small but later receives non-negligible learned
weight on some edges.  Such an analysis would connect the empirical DABP
advantage directly to the feedback-loop mechanism studied in this paper.
\subsection{Using the Split-only Oscillation}
The theoretical analysis predicts that splitting without damping may create a
small oscillatory attractor rather than convergence.  Although this prevents
MS-SCFG from being a good final solver, the two assignments participating in
the oscillation may still contain useful information.  Table~\ref{tab:merge}
evaluates this by comparing the raw MS-SCFG final cost with two procedures
that merge the two assignments observed at iterations 198 and 199.
\begin{table*}[t]
\centering
\caption{Using the two assignments produced by the split-only oscillation.
Lower is better.  Both merges improve significantly on raw MS-SCFG on every
benchmark (Wilcoxon $p < 10^{-8}$).}
\label{tab:merge}
\begin{tabular}{lrrr}
\toprule
\textbf{Problem}
& \textbf{MS-SCFG}
& \textbf{MGM merge}
& \textbf{Optimal merge} \\
\midrule
Graph coloring
& $668.61$
& $98.19$
& $83.19$ \\
Meeting scheduling
& $81.80$
& $20.98$
& $20.14$ \\
Random dense
& $107330.16$
& $101561.58$
& $101561.58$ \\
Random sparse
& $17288.43$
& $15288.05$
& $15243.23$ \\
Scale-free
& $19703.28$
& $17701.06$
& $17663.82$ \\
\bottomrule
\end{tabular}
\end{table*}
The large improvements in Table~\ref{tab:merge} show that the split-only
oscillation is not arbitrary noise.  It often identifies a small set of candidate
values that contains useful information.  The poor raw performance of MS-SCFG
is therefore explained by incoherent mixing of the two oscillation branches
across variables.  Once these branches are merged coherently, solution quality
improves sharply.  This supports the central claim: splitting creates strong
local decision pressure, but damping or an additional coordination step is needed
to transform that pressure into a stable global assignment.
\subsection{Measuring the Mechanism}
\label{sec:mechanism}
The theory of Sections~\ref{sec:setting} and~\ref{sec:tworoutes} predicts
specific, measurable behavior on benchmark instances, which we instrumented
directly (full protocol, tables, and raw results in
Appendix~\ref{app:tworoutes}).  (i)~\emph{Commitment}: the fraction of
messages with a stable active minimizer rises within ten iterations to
$0.73$--$0.97$ on the split random families (versus $0.10$--$0.36$ unsplit
in the tail, family-dependent), while graph coloring commits only in the
weaker selection-content sense ($0.03$ at iteration $9$, $0.36$ in the
tail) --- commitment is real but partial, which is why the reduction of
Section~\ref{sec:tworoutes} is stated conditionally.  (ii)~\emph{Selection
rule}: on an audited subset of $32$ period-2 runs, the two alternating
solutions mapped to one another under the synchronous local selection rule
exactly, and the rule predicted the next assignment correctly on all
sampled transitions.  (iii)~\emph{Attractor structure}: on the same $32$
runs, the limit pair was a layer-wise local minimum of the alternation
cost $\mathrm{cost}_2$ with zero violating single-variable moves, and on
oscillating bipartite instances the rephasings of
Corollary~\ref{cor:bipartite} were coordinatewise local minima of the
original problem in every case.  (iv)~\emph{Damping}: the smallest grid
$\lambda$ from which all larger values stabilized the assignment tail was
at most $0.4$ on every scanned benchmark instance --- far below the
customary $0.9$, and consistent with Theorem~\ref{thm:lambdastar} ---
while Proposition~\ref{prop:kn} exhibits a family where the required
damping exceeds $0.9$ and grows with degree.
\subsection{Ternary Benchmarks}
We also include a ternary benchmark suite in the appendix.  In these benchmarks
the factors are genuine arity-3 constraints rather than collections of binary
constraints.  The ternary suite is important because the cleanest theoretical
results in this paper are for binary domains and binary local minimizer maps.
For domain size 2, preservation of an active minimizer can force a full neutral
step in which the incoming message difference is cancelled.  With three values,
this clean cancellation need not occur for all rows or columns simultaneously.
The ternary experiments therefore test whether the qualitative split-damped
phenomenon persists when the local minimizer structure is richer and only
semi-neutral behavior may remain.
The ternary plots should not be interpreted as a replacement for the binary
theory.  Rather, they are evidence that the mechanism is not limited to a
single pairwise artifact.  If split-damped variants also stabilize on true
arity-3 factors, then the threshold and minimizer-switching explanation is
likely to extend beyond the simplest binary case, although the full theorem
will require additional case distinctions.
\subsection{Discussion}
The empirical results support the mechanism proposed in the theoretical
sections.  Standard Min-sum performs poorly on cyclic benchmark graphs.
Damping alone improves the result but does not fully exploit the structural
change created by splitting.  Splitting alone changes the message dynamics but
often produces oscillatory behavior.  Combining splitting with damping produces
much lower final costs, especially on structured benchmarks such as graph
coloring and meeting scheduling.
The delayed-splitting results show that splitting is not merely a preprocessing
transformation.  The point at which the split is introduced affects the final
result, suggesting that DMS can first form useful beliefs on the original graph
and that splitting can then accelerate stabilization by changing the local
threshold structure of the factor-node updates.  We note, however, that the
best $K$ varies across benchmarks and that the improvement over immediate
splitting reaches statistical significance only on meeting scheduling and
random dense graphs; selecting $K$ automatically is an open question.
The DABP comparison further strengthens the main claim.  Even learned dynamic
damping and learned neighbor weights do not remove the need for splitting.  The
no-split attentive ablation fails to reach the split-based performance on the
structured benchmarks, and the original DABP design itself uses a split
representation.  The symmetric-split DABP ablation shows that the exact split
ratio also matters on the unstructured families, while on the structured
benchmarks the two ratios are statistically indistinguishable and the symmetric
split often runs faster, providing a cleaner mechanistic baseline.
Understanding this tradeoff will require direct analysis of the learned edge
weights and the messages flowing through the two split copies.
\paragraph{Limitations.}
The formal results of Section~\ref{sec:setting} concern a single split binary
function-node with two-value domains; undamped updates, except for the
conditional Proposition~\ref{prop:damped}; and influence from the remainder of
the graph modeled as a bounded unary perturbation.
Proposition~\ref{prop:asym} extends the absorption analysis to asymmetric
splits, but the unconditional convergence theorem
(Theorem~\ref{thm:convergence}) is proved only for the symmetric case.
Whether the bounded-perturbation condition of Theorem~\ref{thm:dyncvg} holds
on benchmark instances is an empirical question, which the instrumentation of
Section~\ref{sec:mechanism} answers: commitment is partial and
family-dependent.  Accordingly, the oscillation analysis of
Sections~\ref{sec:tworoutes}--\ref{sec:whydamping} is conditional: the
reduction to synchronous local selection assumes full commitment; the
two-solution theorem assumes the selection recursion holds along the run
(verified on every audited instance, but not proved to occur); the damping
threshold concerns fixed, fully committed binary patterns; and
Proposition~\ref{prop:kn} shows that no problem-independent damping value
exists, so the reliability of $\lambda = 0.9$ is a property of the
benchmark families.  Finally, on the
unstructured benchmark families the proposed variants do not match DABP's
final costs; our contribution there is the mechanism explanation and the
distributed-execution compatibility, not state-of-the-art solution quality.

% ========== END sec6_experiments.tex ==========

% ========== BEGIN sec7_appendix.tex ==========
\appendix
\section{Additional Experimental Plots}
\label{app:additional_plots}
\begin{figure*}[t]
\centering
\includegraphics[width=.32\textwidth]{plots/random_sparse_cost.pdf}
\includegraphics[width=.32\textwidth]{plots/random_dense_cost.pdf}
\includegraphics[width=.32\textwidth]{plots/scale_free_cost.pdf}
\caption{Mean solution cost per iteration on the random sparse, random dense,
and scale-free binary benchmarks.  These plots show that the split-damped and
delayed-splitting variants remain competitive beyond the structured graph
coloring and meeting scheduling benchmarks.}
\label{Fig:binary_random_cost_curves}
\end{figure*}
\begin{figure*}[t]
\centering
\includegraphics[width=.32\textwidth]{plots/random_sparse_cost_zoom.pdf}
\includegraphics[width=.32\textwidth]{plots/random_dense_cost_zoom.pdf}
\includegraphics[width=.32\textwidth]{plots/scale_free_cost_zoom.pdf}
\caption{Zoomed lower-cost region for the random sparse, random dense, and
scale-free binary benchmarks.  The zoomed plots highlight the differences
between DMS-SCFG, delayed splitting, DABP, DABP-NoSplit, and DABP-SymSplit
near the end of the run.}
\label{Fig:binary_random_cost_curves_zoom}
\end{figure*}
\begin{figure*}[t]
\centering
\includegraphics[width=.32\textwidth]{plots/random_sparse_ternary_cost.pdf}
\includegraphics[width=.32\textwidth]{plots/random_dense_ternary_cost.pdf}
\includegraphics[width=.32\textwidth]{plots/meeting_scheduling_ternary_cost.pdf}
\caption{Mean solution cost per iteration on true arity-3 benchmarks:
random sparse ternary, random dense ternary, and meeting scheduling ternary.
These experiments test whether the split-damped phenomenon persists beyond
binary pairwise factors.}
\label{Fig:ternary_cost_curves}
\end{figure*}
\begin{figure*}[t]
\centering
\includegraphics[width=.32\textwidth]{plots/random_sparse_ternary_cost_zoom.pdf}
\includegraphics[width=.32\textwidth]{plots/random_dense_ternary_cost_zoom.pdf}
\includegraphics[width=.32\textwidth]{plots/meeting_scheduling_ternary_cost_zoom.pdf}
\caption{Zoomed lower-cost region for the true arity-3 benchmarks.  The
ternary results are important because in domain size 3 the binary neutral-step
mechanism becomes only partial: the active minimizer structure is richer, but
the empirical split-damped stabilization pattern still appears.}
\label{Fig:ternary_cost_curves_zoom}
\end{figure*}

% ========== END sec7_appendix.tex ==========

% ========== BEGIN secA_oscillation_appendix.tex ==========
% ======================================================================
% Appendix: proofs and measurement details for the two-solution theory
% (Sections \ref{sec:tworoutes}-\ref{sec:whydamping}).  Assumes
% \appendix was already issued (sec7_appendix.tex).
% ======================================================================
\section{Proofs and Measurement Details for the Two-Solution Analysis}
\label{app:tworoutes}

\subsection{Deferred Proofs}

\begin{proof}[Proof of Lemma~\ref{lem:clones}]
Induction on $t$.  At initialization all messages are zero.  If all
corresponding messages are equal at $t-1$, the two copies receive equal
inputs and hold equal tables, so their $R^{t}$ are equal; the two $Q^{t}$
from $X_i$ toward the pair differ only by which copy's $R^{t-1}$ is
excluded, and those are equal by hypothesis.
\end{proof}

\begin{proof}[Proof of Lemma~\ref{lem:decode}]
Each copy forwards half of the same row:
$R^{t}_{F'_{ij}\to X_i}(v) = \tfrac12 C_{ij}(v, u^{*t}_j) +
\text{const}$, and likewise for $F''_{ij}$ with the same minimizer
(Lemma~\ref{lem:clones}).  Summing the two copies over all neighbors
reconstructs $\sum_j C_{ij}(v, u^{*t}_j)$ up to an additive constant,
which does not affect the argmin.
\end{proof}

\begin{proof}[Proof of Lemma~\ref{lem:ema}]
Solving $m_{\mathrm{even}} = \lambda m_{\mathrm{odd}} + (1-\lambda)u$ and
$m_{\mathrm{odd}} = \lambda m_{\mathrm{even}} + (1-\lambda)v$ gives
$m_{\mathrm{even}} = (u+\lambda v)/(1+\lambda)$ and
$m_{\mathrm{odd}} = (v+\lambda u)/(1+\lambda)$; the deviation from this
alternating pair contracts by $\lambda$ per step.
\end{proof}

\begin{proof}[Proof of Theorem~\ref{thm:lambdastar}]
Within the pattern the computed difference of each flipping message
alternates between the constants $\Delta_1, \Delta_2$, so by
Lemma~\ref{lem:ema} its damped difference converges to the alternating
pair $(\Delta_1 + \lambda\Delta_2)/(1+\lambda)$,
$(\Delta_2 + \lambda\Delta_1)/(1+\lambda)$.  Pattern consistency requires
the first in the upper regime and the second in the lower:
$\Delta_1 + \lambda\Delta_2 \ge (1+\lambda)\tau_U$ and
$\Delta_2 + \lambda\Delta_1 \le (1+\lambda)\tau_L$, i.e., the stated
bounds.  Under strict inequalities the pair sits strictly inside the
regimes, drives are locally constant, and each difference follows an
affine contraction of rate $\lambda$; hence the orbit exists and attracts
a neighborhood.  Fixed-point invariance:
$m = \lambda m + (1-\lambda)\widehat{m}$ iff $m = \widehat{m}$.
\end{proof}

\begin{proof}[Proof of Proposition~\ref{prop:kn}]
By symmetry of the instance and zero initialization (the induction of
Lemma~\ref{lem:clones} extended over the automorphisms of the complete
graph), all $Q$-differences are equal at every iteration; call the
common value $q_t$.  Each split table is $5$ on the diagonal and $0$ off
it, so its difference map is $-\operatorname{clip}(q, -5, 5)$ with
transition interval $[-5, 5]$, and a $Q$-message sums the unary gap and
the $R$-differences of $2(n-1) - 1 = 2n-3$ incoming copies, giving the
stated scalar recurrence.  While $|q_t| \ge 5$ the computed drive
alternates between $\Delta_1 = 0.01 + 5(2n-3)$ and
$\Delta_2 = 0.01 - 5(2n-3)$; Theorem~\ref{thm:lambdastar} with
$\tau_U = -\tau_L = 5$ yields the threshold
$(10(n-2) - 0.01)/(10(n-1) + 0.01)$ after simplification, and local
attraction with two-step contraction $\lambda^{2}$.  Reachability from
zero initialization: inside the strip the recurrence is affine with slope
$\lambda - (1-\lambda)(2n-3)$, which is $\le -1$ exactly when
$\lambda \le (n-2)/(n-1)$; below the threshold the interior fixed point
is repelling and the oscillating iterates leave the strip in finitely
many steps, after which the committed alternation attracts.  At
$\lambda = 0.95$ (for $n = 12$) the interior slope is $-0.10$ and the
message differences converge to the interior fixed point, whose decoded
beliefs are exactly tied --- an artifact of the perfectly uniform
instance --- so assignments there are determined by tie-breaking rather
than by a surviving alternation.
\end{proof}

\begin{proof}[Proof of Corollary~\ref{cor:bipartite}]
For an edge with $i \in A$, $j \in B$: $x'_i = x_i$, $x'_j = y_j$,
$y'_i = y_i$, $y'_j = x_j$, so the two cross terms
$C_{ij}(x_i, y_j) + C_{ij}(y_i, x_j)$ are literally the two intra terms
$C_{ij}(x'_i, x'_j) + C_{ij}(y'_i, y'_j)$; unary terms match because
each variable contributes both of its values once on either side.  A
single-variable move in $x'$ corresponds to a single-layer move in
$(x,y)$ --- in layer $x$ if the variable is in $A$, in layer $y$
otherwise --- so layer-wise local minimality of $\mathrm{cost}_2$
transfers to coordinatewise local minimality of $\mathrm{cost}$ for $x'$
and $y'$ separately.
\end{proof}

\subsection{Measurement Protocol}

All measurements were produced with a vectorized reimplementation of the
synchronous two-phase schedule, verified to produce assignment
trajectories identical to the repository engines
(\texttt{BPEngine}, \texttt{SplitEngine}, \texttt{DampingEngine},
\texttt{DampingSCFGEngine}) on shared instances, with per-iteration
message normalization (without which undamped raw messages lose
floating-point precision by iteration ${\sim}30$).  Families: random
Erd\H{o}s--R\'enyi with $|D_i| = 10$ and uniform integer costs in
$[100, 200)$ at densities $0.25$ (sparse) and $0.6$ (dense); random
binary ($|D_i| = 2$) at density $0.25$; 3-coloring at density $0.15$
with conflict cost $10$; tie-breaking unary preferences uniform in
$[0, 10^{-2})$.  \emph{Period census}: $700$ iterations; the period of
the selected-assignment sequence is the smallest $p \le 64$ that repeats
over the final $150$ iterations (``detected periodic tail''); $12$--$20$
seeds per cell over $n \in \{10, 20, 40, 80\}$.  \emph{Commitment}: the
fraction of messages with a stable active minimizer, at iteration $9$,
$49$, and averaged over the final $50$ iterations, recorded for split
and unsplit runs of the same instances.  \emph{Audited mechanism
subset}: $34$ instances ($10$ sparse, $10$ coloring, $8$ binary, $6$
dense), of which $32$ had period-2 tails; on each we checked the
selection-rule closure on the limit pair (both directions, exact), the
selection-rule prediction on every seventh transition of the recorded
$100$-iteration tail, and layer-wise local minimality of
$\mathrm{cost}_2$ under all single-variable moves.  \emph{Damping scans}:
$19$ instances ($17$ oscillating undamped), grid
$\lambda \in \{0.02, 0.05, 0.08, 0.12, 0.2, 0.3, 0.4, 0.5, 0.6, 0.7,
0.8, 0.9\}$, $2000$ iterations, ``stable'' meaning the assignment is
constant over the final $300$.  \emph{Bipartite verification}: $40$
random complete-bipartite instances ($4{+}4$ variables, $|D_i| = 3$,
uniform integer costs in $[0,100)$), of which $8$ oscillated; on all
$8$, the identity of Corollary~\ref{cor:bipartite} held exactly and both
rephasings were coordinatewise local minima of the original problem
(raw snapshot costs $532$--$1044$ versus rephased costs $474$--$560$).
Raw results and scripts accompany the code
(\texttt{analysis\_oscillation/}).

\subsection{Period Census}

\begin{table}[t]
\centering
\caption{Fraction of runs whose selected-assignment sequence has a
detected periodic tail of period 1 / period 2 / longer / none (window
$150$, $p \le 64$).  MS-SCFG uses the symmetric split, no damping.}
\label{tab:census}
\scriptsize
\begin{tabular}{llrrrrrrrr}
\toprule
& & \multicolumn{4}{c}{MS-SCFG} & \multicolumn{4}{c}{MS} \\
\cmidrule(lr){3-6}\cmidrule(lr){7-10}
family & $n$ & P1 & P2 & P$>$2 & none & P1 & P2 & P$>$2 & none \\
\midrule
sparse dom-10 & 10 & .35 & .65 & 0 & 0 & .65 & 0 & .15 & .20 \\
 & 20 & 0 & 1.00 & 0 & 0 & .10 & .05 & 0 & .85 \\
 & 40 & 0 & 1.00 & 0 & 0 & 0 & 0 & 0 & 1.00 \\
 & 80 & 0 & 1.00 & 0 & 0 & 0 & 0 & 0 & 1.00 \\
dense dom-10 & 10 & .08 & .83 & .08 & 0 & .17 & .17 & 0 & .67 \\
 & 20 & .17 & .83 & 0 & 0 & .08 & 0 & 0 & .92 \\
 & 40 & 0 & 1.00 & 0 & 0 & 0 & 0 & 0 & 1.00 \\
coloring & 10 & .55 & .45 & 0 & 0 & .90 & .05 & .05 & 0 \\
 & 20 & .15 & .85 & 0 & 0 & .45 & .15 & .25 & .15 \\
 & 40 & 0 & 1.00 & 0 & 0 & 0 & .95 & 0 & .05 \\
 & 80 & 0 & 1.00 & 0 & 0 & 0 & 1.00 & 0 & 0 \\
binary & 10 & .85 & .15 & 0 & 0 & .80 & 0 & .20 & 0 \\
 & 20 & .95 & .05 & 0 & 0 & .95 & 0 & 0 & .05 \\
 & 40 & .80 & .15 & .05 & 0 & .65 & .15 & 0 & .20 \\
 & 80 & .35 & .65 & 0 & 0 & .20 & .20 & .10 & .50 \\
\bottomrule
\end{tabular}
\end{table}

Table~\ref{tab:census} gives the full census.  Pooled over all cells,
$191$ of $276$ split runs ($69\%$) had a detected period-2 tail, $83$
period-1, and two longer periods ($16$, $42$); the split runs never
wandered.  Restricted to the $n \ge 40$ cells of the domain-10 and
coloring families the period-2 count is $92/92$.  Unsplit coloring at
$n \ge 40$ also locks into period 2 --- its bounded conflict penalty
lets high-degree fields commit messages without splitting, consistent
with the theory (whenever commitment happens, the period is $1$ or
$2$); unsplit domain-10 runs almost always show no detected period.

\subsection{Damping Grid Audit}

Per-instance smallest grid $\lambda$ from which all larger tested values
were stable: random sparse $\{.12, .2, .2, .2, .2\}$; random dense
$\{.2, .2, .4\}$; coloring $\{.02, .3, .4, .4, .4\}$; binary
$\{.02, .08, .12, .2\}$.  Two dense instances were non-monotone (stable
at one grid value, oscillating at a larger one, stable thereafter),
which is why the ``all larger values'' criterion is used.  On instances
where the number of flipping variables was unchanged from the undamped
run --- a proxy for an unchanged pattern, not a verification of it ---
the residual swing of the damped $Q$-differences followed the
$(1-\lambda)/(1+\lambda)$ law of Lemma~\ref{lem:ema} within a mean
relative deviation of about $4\%$ (random families; coloring
re-equilibrates into different patterns at any tested $\lambda > 0$ and
the same-pattern premise fails there).  Greedy per-variable rephasing of
the limit pair (minimizing the sum of the two layers' ordinary costs)
improved that sum by, on average, $2427$ (sparse), $4096$ (dense),
$374$ (coloring), and $412$ (binary), flipping $6$--$36\%$ of the
variables.

% ========== END secA_oscillation_appendix.tex ==========

\bibliography{refs,DisCSP_refs,DampIns_AIJ,oscillation_refs}
% Reproducibility Checklist (AAAI kit; do not modify the file itself).
% TODO(authors): every answer in ReproducibilityChecklist.tex is still
% "Type your response here" — ALL must be filled before submission.
% With the current draft the honest answers are mostly yes/partial;
% 4.10-4.12 are now "yes" (50 seeds, std devs, paired t + Wilcoxon).
\def\isChecklistMainFile{}
% \input{ReproducibilityChecklist} % (kit file; keep as \input in the split build)
\end{document}
